% Options for packages loaded elsewhere
\PassOptionsToPackage{unicode}{hyperref}
\PassOptionsToPackage{hyphens}{url}
\PassOptionsToPackage{dvipsnames,svgnames*,x11names*}{xcolor}
%
\documentclass[
  a4paper,
  UTF8]{ctexart}
\usepackage{amsmath,amssymb}
\usepackage{lmodern}
\usepackage{iftex}
\ifPDFTeX
  \usepackage[T1]{fontenc}
  \usepackage[utf8]{inputenc}
  \usepackage{textcomp} % provide euro and other symbols
\else % if luatex or xetex
  \usepackage{unicode-math}
  \defaultfontfeatures{Scale=MatchLowercase}
  \defaultfontfeatures[\rmfamily]{Ligatures=TeX,Scale=1}
\fi
% Use upquote if available, for straight quotes in verbatim environments
\IfFileExists{upquote.sty}{\usepackage{upquote}}{}
\IfFileExists{microtype.sty}{% use microtype if available
  \usepackage[]{microtype}
  \UseMicrotypeSet[protrusion]{basicmath} % disable protrusion for tt fonts
}{}
\makeatletter
\@ifundefined{KOMAClassName}{% if non-KOMA class
  \IfFileExists{parskip.sty}{%
    \usepackage{parskip}
  }{% else
    \setlength{\parindent}{0pt}
    \setlength{\parskip}{6pt plus 2pt minus 1pt}}
}{% if KOMA class
  \KOMAoptions{parskip=half}}
\makeatother
\usepackage{xcolor}
\IfFileExists{xurl.sty}{\usepackage{xurl}}{} % add URL line breaks if available
\IfFileExists{bookmark.sty}{\usepackage{bookmark}}{\usepackage{hyperref}}
\hypersetup{
  pdftitle={详细设计说明书},
  pdfauthor={DBMS 控制台项目},
  colorlinks=true,
  linkcolor={blue},
  filecolor={Maroon},
  citecolor={Blue},
  urlcolor={blue},
  pdfcreator={LaTeX via pandoc}}
\urlstyle{same} % disable monospaced font for URLs
\usepackage[margin=2.5cm]{geometry}
\usepackage{color}
\usepackage{fancyvrb}
\newcommand{\VerbBar}{|}
\newcommand{\VERB}{\Verb[commandchars=\\\{\}]}
\DefineVerbatimEnvironment{Highlighting}{Verbatim}{commandchars=\\\{\}}
% Add ',fontsize=\small' for more characters per line
\newenvironment{Shaded}{}{}
\newcommand{\AlertTok}[1]{\textcolor[rgb]{1.00,0.00,0.00}{\textbf{#1}}}
\newcommand{\AnnotationTok}[1]{\textcolor[rgb]{0.38,0.63,0.69}{\textbf{\textit{#1}}}}
\newcommand{\AttributeTok}[1]{\textcolor[rgb]{0.49,0.56,0.16}{#1}}
\newcommand{\BaseNTok}[1]{\textcolor[rgb]{0.25,0.63,0.44}{#1}}
\newcommand{\BuiltInTok}[1]{#1}
\newcommand{\CharTok}[1]{\textcolor[rgb]{0.25,0.44,0.63}{#1}}
\newcommand{\CommentTok}[1]{\textcolor[rgb]{0.38,0.63,0.69}{\textit{#1}}}
\newcommand{\CommentVarTok}[1]{\textcolor[rgb]{0.38,0.63,0.69}{\textbf{\textit{#1}}}}
\newcommand{\ConstantTok}[1]{\textcolor[rgb]{0.53,0.00,0.00}{#1}}
\newcommand{\ControlFlowTok}[1]{\textcolor[rgb]{0.00,0.44,0.13}{\textbf{#1}}}
\newcommand{\DataTypeTok}[1]{\textcolor[rgb]{0.56,0.13,0.00}{#1}}
\newcommand{\DecValTok}[1]{\textcolor[rgb]{0.25,0.63,0.44}{#1}}
\newcommand{\DocumentationTok}[1]{\textcolor[rgb]{0.73,0.13,0.13}{\textit{#1}}}
\newcommand{\ErrorTok}[1]{\textcolor[rgb]{1.00,0.00,0.00}{\textbf{#1}}}
\newcommand{\ExtensionTok}[1]{#1}
\newcommand{\FloatTok}[1]{\textcolor[rgb]{0.25,0.63,0.44}{#1}}
\newcommand{\FunctionTok}[1]{\textcolor[rgb]{0.02,0.16,0.49}{#1}}
\newcommand{\ImportTok}[1]{#1}
\newcommand{\InformationTok}[1]{\textcolor[rgb]{0.38,0.63,0.69}{\textbf{\textit{#1}}}}
\newcommand{\KeywordTok}[1]{\textcolor[rgb]{0.00,0.44,0.13}{\textbf{#1}}}
\newcommand{\NormalTok}[1]{#1}
\newcommand{\OperatorTok}[1]{\textcolor[rgb]{0.40,0.40,0.40}{#1}}
\newcommand{\OtherTok}[1]{\textcolor[rgb]{0.00,0.44,0.13}{#1}}
\newcommand{\PreprocessorTok}[1]{\textcolor[rgb]{0.74,0.48,0.00}{#1}}
\newcommand{\RegionMarkerTok}[1]{#1}
\newcommand{\SpecialCharTok}[1]{\textcolor[rgb]{0.25,0.44,0.63}{#1}}
\newcommand{\SpecialStringTok}[1]{\textcolor[rgb]{0.73,0.40,0.53}{#1}}
\newcommand{\StringTok}[1]{\textcolor[rgb]{0.25,0.44,0.63}{#1}}
\newcommand{\VariableTok}[1]{\textcolor[rgb]{0.10,0.09,0.49}{#1}}
\newcommand{\VerbatimStringTok}[1]{\textcolor[rgb]{0.25,0.44,0.63}{#1}}
\newcommand{\WarningTok}[1]{\textcolor[rgb]{0.38,0.63,0.69}{\textbf{\textit{#1}}}}
\usepackage{longtable,booktabs,array}
\usepackage{calc} % for calculating minipage widths
% Correct order of tables after \paragraph or \subparagraph
\usepackage{etoolbox}
\makeatletter
\patchcmd\longtable{\par}{\if@noskipsec\mbox{}\fi\par}{}{}
\makeatother
% Allow footnotes in longtable head/foot
\IfFileExists{footnotehyper.sty}{\usepackage{footnotehyper}}{\usepackage{footnote}}
\makesavenoteenv{longtable}
\setlength{\emergencystretch}{3em} % prevent overfull lines
\providecommand{\tightlist}{%
  \setlength{\itemsep}{0pt}\setlength{\parskip}{0pt}}
\setcounter{secnumdepth}{5}
\ifLuaTeX
  \usepackage{selnolig}  % disable illegal ligatures
\fi

\title{详细设计说明书}
\author{DBMS 控制台项目}
\date{}

\begin{document}
\maketitle

\renewcommand*\contentsname{目录}
{
\hypersetup{linkcolor=}
\setcounter{tocdepth}{3}
\tableofcontents
}
\hypertarget{ux8be6ux7ec6ux8bbeux8ba1}{%
\section{详细设计}\label{ux8be6ux7ec6ux8bbeux8ba1}}

\hypertarget{ux6587ux6863ux5b9aux4f4d}{%
\subsection{1. 文档定位}\label{ux6587ux6863ux5b9aux4f4d}}

本文档承接以下文档：

\begin{itemize}
\tightlist
\item
  \texttt{需求规格说明书.md}：说明系统要实现什么。
\item
  \texttt{概要设计.md}：说明系统总体怎么拆分。
\item
  \texttt{数据库设计.md}：说明系统数据如何持久化。
\item
  \texttt{接口规范.md}：说明前后端如何通信。
\item
  \texttt{安全设计.md}：说明认证、授权、SQL、文件、容量限制等安全规则。
\end{itemize}

本文档重点说明``如何落地实现''，包括模块职责、关键数据结构、核心流程、伪代码、状态设计和实现顺序。

\hypertarget{ux8be6ux7ec6ux8bbeux8ba1ux76eeux6807}{%
\subsection{2.
详细设计目标}\label{ux8be6ux7ec6ux8bbeux8ba1ux76eeux6807}}

\hypertarget{dd-001-ux5b9eux73b0ux53efux8ffdux8e2a}{%
\subsubsection{DD-001
实现可追踪}\label{dd-001-ux5b9eux73b0ux53efux8ffdux8e2a}}

每个核心模块都应能追踪到需求编号、接口编号、安全编号或数据设计编号。

\hypertarget{dd-002-ux5b9eux73b0ux53efux5206ux9636ux6bb5}{%
\subsubsection{DD-002
实现可分阶段}\label{dd-002-ux5b9eux73b0ux53efux5206ux9636ux6bb5}}

设计必须支持从最小闭环逐步扩展：

\begin{Shaded}
\begin{Highlighting}[]
\NormalTok{认证与工作台}
\NormalTok{  {-}\textgreater{} 数据库管理}
\NormalTok{  {-}\textgreater{} 表结构管理}
\NormalTok{  {-}\textgreater{} 表数据 CRUD}
\NormalTok{  {-}\textgreater{} 查询构造器}
\NormalTok{  {-}\textgreater{} 事务}
\NormalTok{  {-}\textgreater{} 统计和 DIY 工作台}
\end{Highlighting}
\end{Shaded}

\hypertarget{dd-003-ux5b9eux73b0ux53efux7ef4ux62a4}{%
\subsubsection{DD-003
实现可维护}\label{dd-003-ux5b9eux73b0ux53efux7ef4ux62a4}}

后端 controller 不直接拼
SQL，不直接写复杂业务；前端页面不直接堆全部逻辑，复杂逻辑下沉到
composable、store 和组件。

\hypertarget{dd-004-ux5b9eux73b0ux5b89ux5168ux9ed8ux8ba4}{%
\subsubsection{DD-004
实现安全默认}\label{dd-004-ux5b9eux73b0ux5b89ux5168ux9ed8ux8ba4}}

所有数据库、表、字段、视图、索引、约束、资源都必须基于当前用户上下文做归属校验。

\hypertarget{dd-005-ux9759ux6001ux8bbeux8ba1ux4e0eux52a8ux6001ux8bbeux8ba1ux663eux5f0fux5206ux5c42}{%
\subsubsection{DD-005
静态设计与动态设计显式分层}\label{dd-005-ux9759ux6001ux8bbeux8ba1ux4e0eux52a8ux6001ux8bbeux8ba1ux663eux5f0fux5206ux5c42}}

详细设计同时包含静态设计和动态设计。

静态设计描述系统由哪些结构组成，包括：

\begin{longtable}[]{@{}ll@{}}
\toprule
静态设计内容 & 所在位置 \\
\midrule
\endhead
后端目录、controller、service、repository、SQL 工具 &
本文档``后端详细设计'' \\
前端目录、store、composable、layout、component &
本文档``前端详细设计'' \\
系统库表、字段、约束、索引 & \texttt{数据库设计.md} \\
接口路径、请求、响应、错误码 & \texttt{接口规范.md} \\
\bottomrule
\end{longtable}

动态设计描述系统运行时如何流转，包括：

\begin{longtable}[]{@{}ll@{}}
\toprule
动态设计内容 & 所在位置 \\
\midrule
\endhead
请求处理链路 & 本文档``请求处理链路'' \\
登录、刷新 Token、退出登录 & 本文档``认证与会话详细设计'' \\
创建数据库和失败补偿 & 本文档``数据库管理详细设计'' \\
行内编辑草稿、提交、取消 & 本文档``表数据详细设计'' \\
查询 AST 校验、作用域解析、SQL 生成 & 本文档``结构化查询详细设计'' \\
事务开启、提交、回滚、超时回滚 & 本文档``事务详细设计'' \\
前后端交互流程 & 本文档``前后端交互流程'' \\
\bottomrule
\end{longtable}

\hypertarget{ux540eux7aefux8be6ux7ec6ux8bbeux8ba1}{%
\subsection{3.
后端详细设计}\label{ux540eux7aefux8be6ux7ec6ux8bbeux8ba1}}

\hypertarget{ux63a8ux8350ux76eeux5f55ux7ed3ux6784}{%
\subsubsection{3.1
推荐目录结构}\label{ux63a8ux8350ux76eeux5f55ux7ed3ux6784}}

\begin{Shaded}
\begin{Highlighting}[]
\NormalTok{backend/src/}
\NormalTok{  app.ts}
\NormalTok{  server.ts}

\NormalTok{  config/}
\NormalTok{    env.ts}
\NormalTok{    db.ts}
\NormalTok{    limits.ts}

\NormalTok{  middlewares/}
\NormalTok{    authMiddleware.ts}
\NormalTok{    errorMiddleware.ts}
\NormalTok{    traceMiddleware.ts}
\NormalTok{    rateLimitMiddleware.ts}

\NormalTok{  routes/}
\NormalTok{    authRoutes.ts}
\NormalTok{    workbenchRoutes.ts}
\NormalTok{    databaseRoutes.ts}
\NormalTok{    tableRoutes.ts}
\NormalTok{    rowRoutes.ts}
\NormalTok{    queryRoutes.ts}
\NormalTok{    transactionRoutes.ts}
\NormalTok{    statsRoutes.ts}
\NormalTok{    auditRoutes.ts}
\NormalTok{    assetRoutes.ts}

\NormalTok{  controllers/}
\NormalTok{    authController.ts}
\NormalTok{    workbenchController.ts}
\NormalTok{    databaseController.ts}
\NormalTok{    tableController.ts}
\NormalTok{    rowController.ts}
\NormalTok{    queryController.ts}
\NormalTok{    transactionController.ts}
\NormalTok{    statsController.ts}
\NormalTok{    auditController.ts}
\NormalTok{    assetController.ts}

\NormalTok{  services/}
\NormalTok{    authService.ts}
\NormalTok{    tokenService.ts}
\NormalTok{    userService.ts}
\NormalTok{    preferenceService.ts}
\NormalTok{    assetService.ts}
\NormalTok{    databaseService.ts}
\NormalTok{    metadataService.ts}
\NormalTok{    tableService.ts}
\NormalTok{    rowService.ts}
\NormalTok{    queryService.ts}
\NormalTok{    transactionService.ts}
\NormalTok{    statsService.ts}
\NormalTok{    auditService.ts}
\NormalTok{    limitService.ts}

\NormalTok{  repositories/}
\NormalTok{    userRepository.ts}
\NormalTok{    refreshTokenRepository.ts}
\NormalTok{    preferenceRepository.ts}
\NormalTok{    assetRepository.ts}
\NormalTok{    userDatabaseRepository.ts}
\NormalTok{    auditRepository.ts}
\NormalTok{    transactionRepository.ts}

\NormalTok{  sql/}
\NormalTok{    identifier.ts}
\NormalTok{    escape.ts}
\NormalTok{    schemaBuilder.ts}
\NormalTok{    rowBuilder.ts}
\NormalTok{    queryAst.ts}
\NormalTok{    queryBuilder.ts}
\NormalTok{    conditionBuilder.ts}

\NormalTok{  types/}
\NormalTok{    auth.ts}
\NormalTok{    request.ts}
\NormalTok{    response.ts}
\NormalTok{    database.ts}
\NormalTok{    table.ts}
\NormalTok{    row.ts}
\NormalTok{    query.ts}
\NormalTok{    error.ts}

\NormalTok{  utils/}
\NormalTok{    response.ts}
\NormalTok{    errors.ts}
\NormalTok{    crypto.ts}
\NormalTok{    time.ts}
\end{Highlighting}
\end{Shaded}

\hypertarget{ux8bf7ux6c42ux5904ux7406ux94feux8def}{%
\subsubsection{3.2
请求处理链路}\label{ux8bf7ux6c42ux5904ux7406ux94feux8def}}

所有业务请求按固定顺序处理：

\begin{Shaded}
\begin{Highlighting}[]
\NormalTok{HTTP Request}
\NormalTok{  {-}\textgreater{} traceMiddleware 生成 traceId}
\NormalTok{  {-}\textgreater{} rateLimitMiddleware 做基础频率限制}
\NormalTok{  {-}\textgreater{} express.json / multipart 解析}
\NormalTok{  {-}\textgreater{} authMiddleware 解析当前用户}
\NormalTok{  {-}\textgreater{} route}
\NormalTok{  {-}\textgreater{} controller 读取参数}
\NormalTok{  {-}\textgreater{} service 执行业务逻辑}
\NormalTok{  {-}\textgreater{} repository / sql builder / mysql pool}
\NormalTok{  {-}\textgreater{} auditService 记录审计}
\NormalTok{  {-}\textgreater{} response.success 返回}
\NormalTok{  {-}\textgreater{} errorMiddleware 统一错误响应}
\end{Highlighting}
\end{Shaded}

\hypertarget{tracemiddleware}{%
\subsubsection{3.3 traceMiddleware}\label{tracemiddleware}}

职责：

\begin{itemize}
\tightlist
\item
  为每个请求生成 \texttt{traceId}。
\item
  若前端传入 \texttt{X-Trace-Id}，后端可以复用合法 traceId。
\item
  将 \texttt{traceId} 写入请求上下文。
\item
  响应中必须返回 \texttt{traceId}。
\end{itemize}

伪代码：

\begin{Shaded}
\begin{Highlighting}[]
\KeywordTok{function} \FunctionTok{traceMiddleware}\NormalTok{(req}\OperatorTok{,}\NormalTok{ res}\OperatorTok{,}\NormalTok{ next) \{}
\NormalTok{  req}\OperatorTok{.}\AttributeTok{traceId} \OperatorTok{=} \FunctionTok{normalizeTraceId}\NormalTok{(req}\OperatorTok{.}\AttributeTok{headers}\NormalTok{[}\StringTok{\textquotesingle{}x{-}trace{-}id\textquotesingle{}}\NormalTok{]) }\OperatorTok{??} \FunctionTok{generateTraceId}\NormalTok{()}\OperatorTok{;}
\NormalTok{  res}\OperatorTok{.}\FunctionTok{setHeader}\NormalTok{(}\StringTok{\textquotesingle{}X{-}Trace{-}Id\textquotesingle{}}\OperatorTok{,}\NormalTok{ req}\OperatorTok{.}\AttributeTok{traceId}\NormalTok{)}\OperatorTok{;}
  \FunctionTok{next}\NormalTok{()}\OperatorTok{;}
\NormalTok{\}}
\end{Highlighting}
\end{Shaded}

\hypertarget{authmiddleware}{%
\subsubsection{3.4 authMiddleware}\label{authmiddleware}}

职责：

\begin{itemize}
\tightlist
\item
  从
  \texttt{Authorization:\ Bearer\ \textless{}accessToken\textgreater{}}
  读取 Access Token。
\item
  校验签名、有效期和 payload。
\item
  查询用户状态是否可用。
\item
  将 \texttt{currentUser} 写入请求上下文。
\end{itemize}

请求上下文建议：

\begin{Shaded}
\begin{Highlighting}[]
\KeywordTok{type}\NormalTok{ CurrentUser }\OperatorTok{=}\NormalTok{ \{}
\NormalTok{  userID}\OperatorTok{:} \DataTypeTok{number}\OperatorTok{;}
\NormalTok{  accountNo}\OperatorTok{:} \DataTypeTok{string}\OperatorTok{;}
\NormalTok{  status}\OperatorTok{:} \StringTok{\textquotesingle{}active\textquotesingle{}} \OperatorTok{|} \StringTok{\textquotesingle{}disabled\textquotesingle{}}\OperatorTok{;}
\NormalTok{\}}\OperatorTok{;}
\end{Highlighting}
\end{Shaded}

认证失败返回：

\begin{Shaded}
\begin{Highlighting}[]
\NormalTok{HTTP 401}
\NormalTok{error.code = AUTH\_TOKEN\_MISSING / AUTH\_TOKEN\_EXPIRED / AUTH\_TOKEN\_INVALID}
\end{Highlighting}
\end{Shaded}

\hypertarget{errormiddleware}{%
\subsubsection{3.5 errorMiddleware}\label{errormiddleware}}

职责：

\begin{itemize}
\tightlist
\item
  捕获 controller、service、repository 抛出的错误。
\item
  将系统错误转换为 \texttt{接口规范.md} 的统一错误响应。
\item
  MySQL 错误需要映射为明确错误码。
\item
  不返回密码、Token、真实物理库名、服务端文件绝对路径。
\end{itemize}

错误类型建议：

\begin{Shaded}
\begin{Highlighting}[]
\NormalTok{class AppError extends Error \{}
\NormalTok{  httpStatus}\OperatorTok{:} \DataTypeTok{number}\OperatorTok{;}
\NormalTok{  type}\OperatorTok{:} \DataTypeTok{string}\OperatorTok{;}
\NormalTok{  code}\OperatorTok{:} \DataTypeTok{string}\OperatorTok{;}
\NormalTok{  details}\OperatorTok{?:} \DataTypeTok{string}\OperatorTok{;}
\NormalTok{  fieldErrors}\OperatorTok{?:}\NormalTok{ FieldError[]}\OperatorTok{;}
\NormalTok{  mysql}\OperatorTok{?:}\NormalTok{ MysqlSafeError }\OperatorTok{|} \DataTypeTok{null}\OperatorTok{;}
\NormalTok{\}}
\end{Highlighting}
\end{Shaded}

\hypertarget{ux914dux7f6eux4e0eux9650ux5236ux8be6ux7ec6ux8bbeux8ba1}{%
\subsection{4.
配置与限制详细设计}\label{ux914dux7f6eux4e0eux9650ux5236ux8be6ux7ec6ux8bbeux8ba1}}

\hypertarget{env.ts}{%
\subsubsection{4.1 env.ts}\label{env.ts}}

\texttt{env.ts} 统一读取环境变量。关键环境变量缺失时服务必须启动失败。

必需项：

\begin{Shaded}
\begin{Highlighting}[]
\NormalTok{DB\_HOST}
\NormalTok{DB\_USER}
\NormalTok{DB\_PASSWORD}
\NormalTok{DB\_NAME}
\NormalTok{JWT\_ACCESS\_SECRET}
\NormalTok{JWT\_REFRESH\_SECRET}
\end{Highlighting}
\end{Shaded}

可选项：

\begin{Shaded}
\begin{Highlighting}[]
\NormalTok{PORT}
\NormalTok{CORS\_ORIGIN}
\NormalTok{UPLOAD\_ROOT}
\NormalTok{COOKIE\_SECURE}
\end{Highlighting}
\end{Shaded}

\hypertarget{limits.ts}{%
\subsubsection{4.2 limits.ts}\label{limits.ts}}

容量限制集中定义：

\begin{Shaded}
\begin{Highlighting}[]
\NormalTok{export const limits }\OperatorTok{=}\NormalTok{ \{}
\NormalTok{  maxDatabasesPerUser}\OperatorTok{:} \DecValTok{10}\OperatorTok{,}
\NormalTok{  maxTablesPerDatabase}\OperatorTok{:} \DecValTok{100}\OperatorTok{,}
\NormalTok{  maxViewsPerDatabase}\OperatorTok{:} \DecValTok{50}\OperatorTok{,}
\NormalTok{  maxColumnsPerTable}\OperatorTok{:} \DecValTok{80}\OperatorTok{,}
\NormalTok{  maxIndexesPerTable}\OperatorTok{:} \DecValTok{16}\OperatorTok{,}
\NormalTok{  maxUserStorageBytes}\OperatorTok{:} \DecValTok{1024} \OperatorTok{*} \DecValTok{1024} \OperatorTok{*} \DecValTok{1024}\OperatorTok{,}
\NormalTok{  maxDatabaseStorageBytes}\OperatorTok{:} \DecValTok{512} \OperatorTok{*} \DecValTok{1024} \OperatorTok{*} \DecValTok{1024}\OperatorTok{,}
\NormalTok{  maxTableStorageBytes}\OperatorTok{:} \DecValTok{256} \OperatorTok{*} \DecValTok{1024} \OperatorTok{*} \DecValTok{1024}\OperatorTok{,}
\NormalTok{  maxPageSize}\OperatorTok{:} \DecValTok{100}\OperatorTok{,}
\NormalTok{  maxQueryRows}\OperatorTok{:} \DecValTok{1000}\OperatorTok{,}
\NormalTok{  maxBatchInsertRows}\OperatorTok{:} \DecValTok{500}\OperatorTok{,}
\NormalTok{  maxBatchWriteAffectedRows}\OperatorTok{:} \DecValTok{200}\OperatorTok{,}
\NormalTok{  maxJoinCount}\OperatorTok{:} \DecValTok{8}\OperatorTok{,}
\NormalTok{  maxSubqueryDepth}\OperatorTok{:} \DecValTok{3}\OperatorTok{,}
\NormalTok{  maxQueryDurationMs}\OperatorTok{:} \DecValTok{10000}\OperatorTok{,}
\NormalTok{  maxQueryAstBytes}\OperatorTok{:} \DecValTok{64} \OperatorTok{*} \DecValTok{1024}\OperatorTok{,}
\NormalTok{  maxAvatarBytes}\OperatorTok{:} \DecValTok{2} \OperatorTok{*} \DecValTok{1024} \OperatorTok{*} \DecValTok{1024}\OperatorTok{,}
\NormalTok{  maxBackgroundBytes}\OperatorTok{:} \DecValTok{8} \OperatorTok{*} \DecValTok{1024} \OperatorTok{*} \DecValTok{1024}\OperatorTok{,}
\NormalTok{  maxUserAssetBytes}\OperatorTok{:} \DecValTok{100} \OperatorTok{*} \DecValTok{1024} \OperatorTok{*} \DecValTok{1024}\OperatorTok{,}
\NormalTok{\}}\OperatorTok{;}
\end{Highlighting}
\end{Shaded}

\hypertarget{limitservice}{%
\subsubsection{4.3 limitService}\label{limitservice}}

职责：

\begin{itemize}
\tightlist
\item
  创建数据库前检查单用户数据库数量。
\item
  创建表前检查单库表数量。
\item
  新增字段前检查字段数量。
\item
  创建索引前检查索引数量。
\item
  执行查询前检查 AST 体积、JOIN 数、子查询深度、pageSize。
\item
  批量写入前检查请求行数。
\item
  批量修改/删除前检查影响行数。
\item
  上传资源前检查单文件大小和用户资源总量。
\end{itemize}

超限错误：

\begin{Shaded}
\begin{Highlighting}[]
\NormalTok{throw }\KeywordTok{new} \FunctionTok{AppError}\NormalTok{(\{}
\NormalTok{  httpStatus}\OperatorTok{:} \DecValTok{413}\OperatorTok{,}
\NormalTok{  type}\OperatorTok{:} \StringTok{\textquotesingle{}LIMIT\_EXCEEDED\textquotesingle{}}\OperatorTok{,}
\NormalTok{  code}\OperatorTok{:} \StringTok{\textquotesingle{}LIMIT\_EXCEEDED\textquotesingle{}}\OperatorTok{,}
\NormalTok{  details}\OperatorTok{:} \StringTok{\textquotesingle{}单用户数据库数量已达到上限：10\textquotesingle{}}\OperatorTok{,}
\NormalTok{\})}\OperatorTok{;}
\end{Highlighting}
\end{Shaded}

\hypertarget{ux8ba4ux8bc1ux4e0eux4f1aux8bddux8be6ux7ec6ux8bbeux8ba1}{%
\subsection{5.
认证与会话详细设计}\label{ux8ba4ux8bc1ux4e0eux4f1aux8bddux8be6ux7ec6ux8bbeux8ba1}}

\hypertarget{authservice}{%
\subsubsection{5.1 authService}\label{authservice}}

职责：

\begin{itemize}
\tightlist
\item
  注册用户。
\item
  登录校验。
\item
  更新最近登录时间。
\item
  调用 tokenService 签发双 Token。
\item
  调用 auditService 记录认证日志。
\end{itemize}

\hypertarget{tokenservice}{%
\subsubsection{5.2 tokenService}\label{tokenservice}}

职责：

\begin{itemize}
\tightlist
\item
  生成 Access Token。
\item
  生成 Refresh Token。
\item
  哈希 Refresh Token 后写入 \texttt{refresh\_tokens}。
\item
  刷新 Token 时执行 Refresh Token Rotation。
\item
  退出登录时撤销 Refresh Token。
\end{itemize}

\hypertarget{ux767bux5f55ux6d41ux7a0b}{%
\subsubsection{5.3 登录流程}\label{ux767bux5f55ux6d41ux7a0b}}

\begin{Shaded}
\begin{Highlighting}[]
\NormalTok{POST /api/v1/auth/login}
\NormalTok{  {-}\textgreater{} 校验 accountNo、password}
\NormalTok{  {-}\textgreater{} 查询 users}
\NormalTok{  {-}\textgreater{} bcrypt.compare}
\NormalTok{  {-}\textgreater{} 生成 accessToken，15 分钟有效}
\NormalTok{  {-}\textgreater{} 生成 refreshToken，7/30 天有效}
\NormalTok{  {-}\textgreater{} refreshToken 哈希写入 refresh\_tokens}
\NormalTok{  {-}\textgreater{} refreshToken 设置 HttpOnly Cookie}
\NormalTok{  {-}\textgreater{} 返回 accessToken 和用户展示信息}
\end{Highlighting}
\end{Shaded}

\hypertarget{refresh-token-rotation}{%
\subsubsection{5.4 Refresh Token
Rotation}\label{refresh-token-rotation}}

刷新流程：

\begin{Shaded}
\begin{Highlighting}[]
\NormalTok{POST /api/v1/auth/refresh}
\NormalTok{  {-}\textgreater{} 从 HttpOnly Cookie 读取 refreshToken}
\NormalTok{  {-}\textgreater{} 哈希后查询 refresh\_tokens}
\NormalTok{  {-}\textgreater{} 校验未过期、未撤销}
\NormalTok{  {-}\textgreater{} 撤销旧 refreshToken}
\NormalTok{  {-}\textgreater{} 新增新 refreshToken}
\NormalTok{  {-}\textgreater{} 返回新 accessToken}
\NormalTok{  {-}\textgreater{} 设置新 refreshToken Cookie}
\end{Highlighting}
\end{Shaded}

异常处理：

\begin{itemize}
\tightlist
\item
  旧 token 不存在：返回 401。
\item
  旧 token 已撤销但再次使用：可能存在 token 盗用，撤销同一 token\_family
  下所有 token。
\item
  用户状态 disabled：返回 403。
\end{itemize}

\hypertarget{ux7528ux6237ux5de5ux4f5cux53f0ux8be6ux7ec6ux8bbeux8ba1}{%
\subsection{6.
用户工作台详细设计}\label{ux7528ux6237ux5de5ux4f5cux53f0ux8be6ux7ec6ux8bbeux8ba1}}

\hypertarget{workbenchservice}{%
\subsubsection{6.1 workbenchService}\label{workbenchservice}}

工作台概览聚合以下数据：

\begin{itemize}
\tightlist
\item
  当前用户资料。
\item
  用户偏好。
\item
  用户资源引用。
\item
  默认 Dashboard 布局。
\item
  最近访问数据库。
\item
  最近操作日志。
\item
  用户级统计摘要。
\item
  数据库入口列表和数据库级摘要卡片。
\end{itemize}

首页职责边界：

\begin{itemize}
\tightlist
\item
  首页只展示用户层、数据库层统计和数据库入口。
\item
  表结构、表数据、查询构造器等数据库内部操作统一进入
  \texttt{DatabaseWorkbench.vue} 后完成。
\item
  \texttt{Home.vue} 不直接承载表设计表单，避免总览页变成具体操作页。
\end{itemize}

\hypertarget{preferenceservice}{%
\subsubsection{6.2 preferenceService}\label{preferenceservice}}

管理：

\begin{itemize}
\tightlist
\item
  主题模式。
\item
  主题色相。
\item
  背景类型。
\item
  背景资源。
\item
  默认分页大小。
\item
  二次确认偏好。
\end{itemize}

偏好更新必须只允许当前用户更新自己的记录。

\hypertarget{assetservice}{%
\subsubsection{6.3 assetService}\label{assetservice}}

资源类型：

\begin{Shaded}
\begin{Highlighting}[]
\NormalTok{avatar}
\NormalTok{background}
\end{Highlighting}
\end{Shaded}

存储原则：

\begin{itemize}
\tightlist
\item
  文件实体保存到 \texttt{UPLOAD\_ROOT}。
\item
  系统库 \texttt{user\_assets} 保存元信息。
\item
  文件名由服务端随机生成。
\item
  读取资源必须校验 \texttt{asset.user\_id\ =\ currentUserId}。
\item
  允许 png、jpeg、webp；SVG 固定拒绝，避免脚本、外链和内容嗅探风险。
\end{itemize}

\hypertarget{ux6570ux636eux5e93ux7ba1ux7406ux8be6ux7ec6ux8bbeux8ba1}{%
\subsection{7.
数据库管理详细设计}\label{ux6570ux636eux5e93ux7ba1ux7406ux8be6ux7ec6ux8bbeux8ba1}}

\hypertarget{databaseservice.createdatabase}{%
\subsubsection{7.1
databaseService.createDatabase}\label{databaseservice.createdatabase}}

输入：

\begin{Shaded}
\begin{Highlighting}[]
\KeywordTok{type}\NormalTok{ CreateDatabaseInput }\OperatorTok{=}\NormalTok{ \{}
\NormalTok{  userID}\OperatorTok{:} \DataTypeTok{number}\OperatorTok{;}
\NormalTok{  displayName}\OperatorTok{:} \DataTypeTok{string}\OperatorTok{;}
\NormalTok{  charset}\OperatorTok{?:} \DataTypeTok{string}\OperatorTok{;}
\NormalTok{  collation}\OperatorTok{?:} \DataTypeTok{string}\OperatorTok{;}
\NormalTok{\}}\OperatorTok{;}
\end{Highlighting}
\end{Shaded}

流程：

\begin{Shaded}
\begin{Highlighting}[]
\NormalTok{1. 校验 displayName 标识符。}
\NormalTok{2. 检查当前用户数据库数量。}
\NormalTok{3. 检查当前用户 displayName 是否重复。}
\NormalTok{4. 生成 schemaName。}
\NormalTok{5. 执行 CREATE DATABASE。}
\NormalTok{6. 写入 user\_databases。}
\NormalTok{7. 写入 audit\_logs。}
\NormalTok{8. 返回 databaseId 和 displayName。}
\end{Highlighting}
\end{Shaded}

当前落地版本补充：

\begin{itemize}
\tightlist
\item
  创建时真实物理库名格式为
  \texttt{u\{userId\}\_\{safeDisplayName\}\_\{randomSuffix\}}。
\item
  如果随机后缀发生极低概率碰撞，最多重试 5 次。
\item
  服务启动会自动检查 \texttt{user\_databases}
  表，避免开发环境漏执行新版初始化脚本。
\end{itemize}

补偿策略：

\begin{Shaded}
\begin{Highlighting}[]
\NormalTok{如果 CREATE DATABASE 成功但 user\_databases 写入失败：}
\NormalTok{  {-}\textgreater{} 尝试 DROP DATABASE schemaName}
\NormalTok{  {-}\textgreater{} 记录失败审计}
\NormalTok{  {-}\textgreater{} 返回系统错误}
\end{Highlighting}
\end{Shaded}

删除数据库采用二次确认：

\begin{enumerate}
\def\labelenumi{\arabic{enumi}.}
\tightlist
\item
  前端要求用户输入完整数据库显示名。
\item
  请求体传入 \texttt{confirmation.confirmed=true} 和
  \texttt{confirmation.confirmText}。
\item
  后端先把状态改为 \texttt{deleting}。
\item
  删除物理 schema 成功后把状态改为 \texttt{deleted}。
\item
  如果物理 schema 删除失败，状态回滚为 \texttt{active}。
\end{enumerate}

\hypertarget{schemaname-ux751fux6210}{%
\subsubsection{7.2 schemaName 生成}\label{schemaname-ux751fux6210}}

规则：

\begin{Shaded}
\begin{Highlighting}[]
\NormalTok{u\{userID\}\_\{displayName\}\_\{suffix\}}
\end{Highlighting}
\end{Shaded}

其中：

\begin{itemize}
\tightlist
\item
  \texttt{displayName} 已经符合标识符规则。
\item
  \texttt{suffix} 使用短随机字符串，防止重名冲突。
\item
  \texttt{schemaName} 不返回前端。
\end{itemize}

\hypertarget{databaseresolver}{%
\subsubsection{7.3 databaseResolver}\label{databaseresolver}}

所有涉及数据库的接口必须先解析：

\begin{Shaded}
\begin{Highlighting}[]
\NormalTok{async }\KeywordTok{function} \FunctionTok{resolveUserDatabase}\NormalTok{(userID}\OperatorTok{:} \DataTypeTok{number}\OperatorTok{,}\NormalTok{ databaseId}\OperatorTok{:} \DataTypeTok{number}\NormalTok{) \{}
\NormalTok{  const db }\OperatorTok{=}\NormalTok{ await userDatabaseRepository}\OperatorTok{.}\FunctionTok{findActiveByIdAndUser}\NormalTok{(databaseId}\OperatorTok{,}\NormalTok{ userID)}\OperatorTok{;}
  \FunctionTok{if}\NormalTok{ (}\OperatorTok{!}\NormalTok{db) throw }\FunctionTok{forbiddenOrNotFound}\NormalTok{()}\OperatorTok{;}
\NormalTok{  return db}\OperatorTok{;}
\NormalTok{\}}
\end{Highlighting}
\end{Shaded}

默认返回 403 还是 404：

\begin{itemize}
\tightlist
\item
  如果资源存在但不属于当前用户，不应暴露存在性，推荐统一返回 404
  或统一业务提示。
\item
  管理类项目为了调试清晰，可以返回 403，但不能包含真实物理库名。
\end{itemize}

\hypertarget{ux5143ux6570ux636eux4e0eux8868ux7ed3ux6784ux8be6ux7ec6ux8bbeux8ba1}{%
\subsection{8.
元数据与表结构详细设计}\label{ux5143ux6570ux636eux4e0eux8868ux7ed3ux6784ux8be6ux7ec6ux8bbeux8ba1}}

\hypertarget{metadataservice}{%
\subsubsection{8.1 metadataService}\label{metadataservice}}

职责：

\begin{itemize}
\tightlist
\item
  查询表、视图列表。
\item
  查询字段列表。
\item
  查询主键。
\item
  查询索引。
\item
  查询约束。
\item
  判断表/视图/字段是否存在。
\end{itemize}

数据来源：

\begin{Shaded}
\begin{Highlighting}[]
\NormalTok{information\_schema.tables}
\NormalTok{information\_schema.columns}
\NormalTok{information\_schema.statistics}
\NormalTok{information\_schema.table\_constraints}
\NormalTok{information\_schema.key\_column\_usage}
\NormalTok{information\_schema.views}
\end{Highlighting}
\end{Shaded}

\hypertarget{identifier-ux5de5ux5177}{%
\subsubsection{8.2 identifier 工具}\label{identifier-ux5de5ux5177}}

禁止业务代码直接拼接用户输入的数据库名、表名、字段名。

\begin{Shaded}
\begin{Highlighting}[]
\KeywordTok{function} \FunctionTok{assertIdentifier}\NormalTok{(value}\OperatorTok{:} \DataTypeTok{string}\OperatorTok{,}\NormalTok{ label}\OperatorTok{:} \DataTypeTok{string}\NormalTok{)}\OperatorTok{:} \DataTypeTok{void}\OperatorTok{;}
\KeywordTok{function} \FunctionTok{quoteIdentifier}\NormalTok{(value}\OperatorTok{:} \DataTypeTok{string}\NormalTok{)}\OperatorTok{:} \DataTypeTok{string}\OperatorTok{;}
\KeywordTok{function} \FunctionTok{quotePath}\NormalTok{(parts}\OperatorTok{:} \DataTypeTok{string}\NormalTok{[])}\OperatorTok{:} \DataTypeTok{string}\OperatorTok{;}
\end{Highlighting}
\end{Shaded}

规则：

\begin{Shaded}
\begin{Highlighting}[]
\NormalTok{\^{}[A{-}Za{-}z][A{-}Za{-}z0{-}9\_]\{0,63\}$}
\end{Highlighting}
\end{Shaded}

\hypertarget{schemabuilder}{%
\subsubsection{8.3 schemaBuilder}\label{schemabuilder}}

职责：

\begin{itemize}
\tightlist
\item
  根据结构化表定义生成 \texttt{CREATE\ TABLE}。
\item
  根据结构化变更生成 \texttt{ALTER\ TABLE}。
\item
  控制字段类型白名单。
\item
  控制约束、索引、外键策略白名单。
\end{itemize}

\texttt{CREATE\ TABLE} 生成流程：

\begin{Shaded}
\begin{Highlighting}[]
\NormalTok{1. 校验 tableName。}
\NormalTok{2. 校验 columns 长度。}
\NormalTok{3. 对每个 column 校验 name、type、length、precision、scale、defaultValue。}
\NormalTok{4. 校验 constraints 引用字段存在。}
\NormalTok{5. 校验 indexes 引用字段存在。}
\NormalTok{6. 校验最多存在一个主键；无主键表允许创建，但数据编辑阶段必须按无主键表策略降级。}
\NormalTok{7. 生成字段定义 SQL。}
\NormalTok{8. 生成约束 SQL。}
\NormalTok{9. 生成索引 SQL。}
\NormalTok{10. 组合 CREATE TABLE。}
\end{Highlighting}
\end{Shaded}

\hypertarget{ux8868ux6570ux636eux8be6ux7ec6ux8bbeux8ba1}{%
\subsection{9.
表数据详细设计}\label{ux8868ux6570ux636eux8be6ux7ec6ux8bbeux8ba1}}

\hypertarget{rowservice.queryrows}{%
\subsubsection{9.1 rowService.queryRows}\label{rowservice.queryrows}}

职责：

\begin{itemize}
\tightlist
\item
  分页预览表数据。
\item
  查询字段结构。
\item
  校验筛选字段。
\item
  限制 limit 和 offset。
\item
  返回 columns、rows、total、hasMore 和有限值 facets。
\end{itemize}

流程：

\begin{Shaded}
\begin{Highlighting}[]
\NormalTok{1. resolveUserDatabase。}
\NormalTok{2. 校验 tableName。}
\NormalTok{3. metadataService 确认表存在。}
\NormalTok{4. metadataService 查询字段和主键。}
\NormalTok{5. 校验 filters 是否来自真实字段。}
\NormalTok{6. 构造 SELECT 和 WHERE。}
\NormalTok{7. 查询 total。}
\NormalTok{8. 查询分页数据。}
\NormalTok{9. 查询有限离散字段候选值。}
\NormalTok{10. 返回。}
\end{Highlighting}
\end{Shaded}

\hypertarget{ux884cux5185ux7f16ux8f91ux8349ux7a3fux673aux5236}{%
\subsubsection{9.2
行内编辑草稿机制}\label{ux884cux5185ux7f16ux8f91ux8349ux7a3fux673aux5236}}

行内编辑由前端负责草稿态，后端只接收明确提交。

前端状态：

\begin{Shaded}
\begin{Highlighting}[]
\KeywordTok{type}\NormalTok{ RowDraftState }\OperatorTok{=}\NormalTok{ \{}
\NormalTok{  rowKey}\OperatorTok{:} \DataTypeTok{string}\OperatorTok{;}
\NormalTok{  original}\OperatorTok{:} \BuiltInTok{Record}\OperatorTok{\textless{}}\DataTypeTok{string}\OperatorTok{,} \DataTypeTok{unknown}\OperatorTok{\textgreater{};}
\NormalTok{  draft}\OperatorTok{:} \BuiltInTok{Record}\OperatorTok{\textless{}}\DataTypeTok{string}\OperatorTok{,} \DataTypeTok{unknown}\OperatorTok{\textgreater{};}
\NormalTok{  changedFields}\OperatorTok{:} \BuiltInTok{Record}\OperatorTok{\textless{}}\DataTypeTok{string}\OperatorTok{,} \DataTypeTok{unknown}\OperatorTok{\textgreater{};}
\NormalTok{  status}\OperatorTok{:} \StringTok{\textquotesingle{}clean\textquotesingle{}} \OperatorTok{|} \StringTok{\textquotesingle{}dirty\textquotesingle{}} \OperatorTok{|} \StringTok{\textquotesingle{}saving\textquotesingle{}} \OperatorTok{|} \StringTok{\textquotesingle{}error\textquotesingle{}}\OperatorTok{;}
\NormalTok{  errorMessage}\OperatorTok{?:} \DataTypeTok{string}\OperatorTok{;}
\NormalTok{\}}\OperatorTok{;}
\end{Highlighting}
\end{Shaded}

交互流程：

\begin{Shaded}
\begin{Highlighting}[]
\NormalTok{用户修改单元格}
\NormalTok{  {-}\textgreater{} 前端更新 draft}
\NormalTok{  {-}\textgreater{} 对比 original 生成 changedFields}
\NormalTok{  {-}\textgreater{} 行状态变为 dirty}
\NormalTok{  {-}\textgreater{} 显示“更新”和“取消”}

\NormalTok{用户点击更新}
\NormalTok{  {-}\textgreater{} 调用 PATCH /rows}
\NormalTok{  {-}\textgreater{} 请求体包含 primaryKey + set}
\NormalTok{  {-}\textgreater{} 成功后 original = draft}
\NormalTok{  {-}\textgreater{} 行状态变为 clean}

\NormalTok{用户点击取消}
\NormalTok{  {-}\textgreater{} draft = original}
\NormalTok{  {-}\textgreater{} changedFields 清空}
\NormalTok{  {-}\textgreater{} 行状态变为 clean}
\end{Highlighting}
\end{Shaded}

\hypertarget{rowservice.updaterows}{%
\subsubsection{9.3 rowService.updateRows}\label{rowservice.updaterows}}

单行更新：

\begin{Shaded}
\begin{Highlighting}[]
\NormalTok{1. resolveUserDatabase。}
\NormalTok{2. 校验 tableName。}
\NormalTok{3. 读取表结构和主键。}
\NormalTok{4. 校验 primaryKey 完整。}
\NormalTok{5. 校验 set 字段存在且不是禁止修改字段。}
\NormalTok{6. 检查单表容量限制。}
\NormalTok{7. 构造 UPDATE ... WHERE pk \textless{}=\textgreater{} ?。}
\NormalTok{8. 执行更新。}
\NormalTok{9. 影响行数必须等于 1。}
\NormalTok{10. 回读最新行数据。}
\NormalTok{11. 写入审计日志。}
\end{Highlighting}
\end{Shaded}

批量更新：

\begin{Shaded}
\begin{Highlighting}[]
\NormalTok{1. 校验 where GUI 条件对象。}
\NormalTok{2. 先执行 SELECT COUNT(*) 预估影响行数。}
\NormalTok{3. 超过阈值时要求 confirmation。}
\NormalTok{4. 构造 UPDATE。}
\NormalTok{5. 执行并审计。}
\end{Highlighting}
\end{Shaded}

\hypertarget{ux65e0ux4e3bux952eux8868ux7b56ux7565}{%
\subsubsection{9.4
无主键表策略}\label{ux65e0ux4e3bux952eux8868ux7b56ux7565}}

第一阶段建议：

\begin{itemize}
\tightlist
\item
  允许浏览无主键表。
\item
  允许新增数据。
\item
  禁止行内编辑。
\item
  禁止单行删除。
\item
  基于主键列表的批量修改/删除必须限制影响行数并进行二次确认；基于条件的批量修改/删除必须通过
  GUI 条件构造器。
\end{itemize}

原因：

\begin{Shaded}
\begin{Highlighting}[]
\NormalTok{没有主键就无法稳定定位“用户正在编辑的那一行”}
\end{Highlighting}
\end{Shaded}

\hypertarget{v5-ux5f53ux524dux5b9eux73b0ux8303ux56f4}{%
\subsubsection{9.5 V5
当前实现范围}\label{v5-ux5f53ux524dux5b9eux73b0ux8303ux56f4}}

V5 已落地接口：

\begin{itemize}
\tightlist
\item
  \texttt{GET\ /api/v1/databases/:databaseId/tables/:tableName/preview}：分页预览、字段筛选、有限值
  facets。
\item
  \texttt{POST\ /api/v1/databases/:databaseId/tables/:tableName/rows}：新增单行和批量新增。
\item
  \texttt{PATCH\ /api/v1/databases/:databaseId/tables/:tableName/rows}：基于主键更新单行和基于主键列表批量修改。
\item
  \texttt{DELETE\ /api/v1/databases/:databaseId/tables/:tableName/rows}：基于主键删除单行和基于主键列表批量删除。
\item
  \texttt{PATCH\ /api/v1/databases/:databaseId/tables/:tableName/schema}：在线结构管理。
\end{itemize}

V5 已落地前端交互：

\begin{itemize}
\tightlist
\item
  工作台主区域展示表格预览。
\item
  表头筛选留空表示全选，按 Enter 或刷新按钮后请求。
\item
  支持表头排序和列显示控制。
\item
  新增行在表格顶部进入草稿态，保存后刷新预览。
\item
  批量新增通过 JSON 数组或 CSV 文件导入弹窗提交；CSV 文件在前端按
  UTF-8/GBK
  解码，按标题行自动匹配字段，无法匹配时允许用户通过下拉框手动配置 CSV
  列到目标字段的映射关系。
\item
  支持从 CSV
  直接创建并导入新表：用户填写表名，字段可由标题行自动生成或手动输入，字段类型和可空等常用设置通过下拉框/勾选项配置；导入数据失败时回滚删除刚创建的新表。
\item
  行内编辑只修改本地草稿，点击``更新''后才提交。
\item
  批量修改通过已选主键列表和字段值弹窗提交。
\item
  批量删除通过已选主键列表和确认弹窗提交。
\item
  无主键表只允许浏览和新增，编辑与删除按钮禁用。
\end{itemize}

已补齐能力：

\begin{itemize}
\tightlist
\item
  后端支持批量新增、批量修改和批量删除，并使用事务保证批次原子性。
\item
  前端支持排序、列显示控制、JSON/CSV
  批量新增、基于主键列表的批量修改和批量删除。
\item
  结构管理支持字段新增、字段修改、字段删除、索引新增/删除和字符串有限取值检查约束。
\end{itemize}

后续增强：

\begin{itemize}
\tightlist
\item
  基于 GUI 条件构造器的条件批量修改/删除。
\item
  级联影响分析和数据写操作审计日志的完整落库。
\end{itemize}

\hypertarget{ux7ed3ux6784ux5316ux67e5ux8be2ux8be6ux7ec6ux8bbeux8ba1}{%
\subsection{10.
结构化查询详细设计}\label{ux7ed3ux6784ux5316ux67e5ux8be2ux8be6ux7ec6ux8bbeux8ba1}}

\hypertarget{ux67e5ux8be2-ast}{%
\subsubsection{10.1 查询 AST}\label{ux67e5ux8be2-ast}}

前端不提交原始 SQL，而是提交查询 AST。

核心结构：

\begin{Shaded}
\begin{Highlighting}[]
\KeywordTok{type}\NormalTok{ SelectQueryAst }\OperatorTok{=}\NormalTok{ \{}
\NormalTok{  from}\OperatorTok{:}\NormalTok{ QuerySource}\OperatorTok{;}
\NormalTok{  joins}\OperatorTok{?:}\NormalTok{ JoinNode[]}\OperatorTok{;}
\NormalTok{  fields}\OperatorTok{:}\NormalTok{ FieldNode[]}\OperatorTok{;}
\NormalTok{  filters}\OperatorTok{?:}\NormalTok{ ConditionNode[]}\OperatorTok{;}
\NormalTok{  groups}\OperatorTok{?:}\NormalTok{ FieldRef[]}\OperatorTok{;}
\NormalTok{  having}\OperatorTok{?:}\NormalTok{ ConditionNode[]}\OperatorTok{;}
\NormalTok{  sorts}\OperatorTok{?:}\NormalTok{ SortNode[]}\OperatorTok{;}
\NormalTok{  page}\OperatorTok{?:} \DataTypeTok{number}\OperatorTok{;}
\NormalTok{  pageSize}\OperatorTok{?:} \DataTypeTok{number}\OperatorTok{;}
\NormalTok{\}}\OperatorTok{;}

\KeywordTok{type}\NormalTok{ QuerySource }\OperatorTok{=}
  \OperatorTok{|}\NormalTok{ \{ type}\OperatorTok{:} \StringTok{\textquotesingle{}table\textquotesingle{}}\OperatorTok{;}\NormalTok{ table}\OperatorTok{:} \DataTypeTok{string}\OperatorTok{;}\NormalTok{ alias}\OperatorTok{?:} \DataTypeTok{string}\NormalTok{ \}}
  \OperatorTok{|}\NormalTok{ \{ type}\OperatorTok{:} \StringTok{\textquotesingle{}view\textquotesingle{}}\OperatorTok{;}\NormalTok{ view}\OperatorTok{:} \DataTypeTok{string}\OperatorTok{;}\NormalTok{ alias}\OperatorTok{?:} \DataTypeTok{string}\NormalTok{ \}}
  \OperatorTok{|}\NormalTok{ \{ type}\OperatorTok{:} \StringTok{\textquotesingle{}subquery\textquotesingle{}}\OperatorTok{;}\NormalTok{ query}\OperatorTok{:}\NormalTok{ SelectQueryAst}\OperatorTok{;}\NormalTok{ alias}\OperatorTok{:} \DataTypeTok{string}\NormalTok{ \}}\OperatorTok{;}
\end{Highlighting}
\end{Shaded}

条件节点：

\begin{Shaded}
\begin{Highlighting}[]
\KeywordTok{type}\NormalTok{ ConditionNode }\OperatorTok{=}\NormalTok{ \{}
\NormalTok{  field}\OperatorTok{?:} \DataTypeTok{string}\OperatorTok{;}
\NormalTok{  rightField}\OperatorTok{?:} \DataTypeTok{string}\OperatorTok{;}
\NormalTok{  operator}\OperatorTok{:} \StringTok{\textquotesingle{}eq\textquotesingle{}} \OperatorTok{|} \StringTok{\textquotesingle{}ne\textquotesingle{}} \OperatorTok{|} \StringTok{\textquotesingle{}gt\textquotesingle{}} \OperatorTok{|} \StringTok{\textquotesingle{}gte\textquotesingle{}} \OperatorTok{|} \StringTok{\textquotesingle{}lt\textquotesingle{}} \OperatorTok{|} \StringTok{\textquotesingle{}lte\textquotesingle{}} \OperatorTok{|}
    \StringTok{\textquotesingle{}like\textquotesingle{}} \OperatorTok{|} \StringTok{\textquotesingle{}in\textquotesingle{}} \OperatorTok{|} \StringTok{\textquotesingle{}notIn\textquotesingle{}} \OperatorTok{|} \StringTok{\textquotesingle{}between\textquotesingle{}} \OperatorTok{|}
    \StringTok{\textquotesingle{}isNull\textquotesingle{}} \OperatorTok{|} \StringTok{\textquotesingle{}isNotNull\textquotesingle{}} \OperatorTok{|} \StringTok{\textquotesingle{}exists\textquotesingle{}} \OperatorTok{|} \StringTok{\textquotesingle{}notExists\textquotesingle{}}\OperatorTok{;}
\NormalTok{  value}\OperatorTok{?:} \DataTypeTok{unknown}\OperatorTok{;}
\NormalTok{  values}\OperatorTok{?:} \DataTypeTok{unknown}\NormalTok{[]}\OperatorTok{;}
\NormalTok{  subquery}\OperatorTok{?:}\NormalTok{ SelectQueryAst}\OperatorTok{;}
\NormalTok{  logic}\OperatorTok{?:} \StringTok{\textquotesingle{}AND\textquotesingle{}} \OperatorTok{|} \StringTok{\textquotesingle{}OR\textquotesingle{}}\OperatorTok{;}
\NormalTok{\}}\OperatorTok{;}
\end{Highlighting}
\end{Shaded}

\hypertarget{queryvalidator}{%
\subsubsection{10.2 queryValidator}\label{queryvalidator}}

职责：

\begin{itemize}
\tightlist
\item
  校验 AST 体积。
\item
  校验主表存在。
\item
  校验 JOIN 来源存在。
\item
  校验子查询深度不超过 3。
\item
  校验总 JOIN 数不超过 8。
\item
  校验字段引用在当前作用域存在。
\item
  校验聚合函数白名单。
\item
  校验操作符白名单。
\item
  校验分页上限。
\end{itemize}

\hypertarget{ux67e5ux8be2ux4f5cux7528ux57df}{%
\subsubsection{10.3 查询作用域}\label{ux67e5ux8be2ux4f5cux7528ux57df}}

查询构造器需要维护作用域：

\begin{Shaded}
\begin{Highlighting}[]
\NormalTok{scope:}
\NormalTok{  alias {-}\textgreater{} source}
\NormalTok{  alias.field {-}\textgreater{} column metadata}
\end{Highlighting}
\end{Shaded}

当遇到子查询：

\begin{Shaded}
\begin{Highlighting}[]
\NormalTok{1. 创建子 scope。}
\NormalTok{2. 注册子查询 from 和 joins。}
\NormalTok{3. 校验子查询 fields 和 filters。}
\NormalTok{4. 子查询输出字段作为派生表字段暴露给父 scope。}
\end{Highlighting}
\end{Shaded}

\hypertarget{querybuilder}{%
\subsubsection{10.4 queryBuilder}\label{querybuilder}}

生成 SQL 流程：

\begin{Shaded}
\begin{Highlighting}[]
\NormalTok{buildSelect(ast, context)}
\NormalTok{  {-}\textgreater{} validate ast}
\NormalTok{  {-}\textgreater{} build FROM}
\NormalTok{  {-}\textgreater{} build JOIN}
\NormalTok{  {-}\textgreater{} build SELECT fields}
\NormalTok{  {-}\textgreater{} build WHERE}
\NormalTok{  {-}\textgreater{} build GROUP BY}
\NormalTok{  {-}\textgreater{} build HAVING}
\NormalTok{  {-}\textgreater{} build ORDER BY}
\NormalTok{  {-}\textgreater{} build LIMIT OFFSET}
\NormalTok{  {-}\textgreater{} 返回 \{ sql, params \}}
\end{Highlighting}
\end{Shaded}

所有字段、表名、别名使用 \texttt{quoteIdentifier}。所有值放入
\texttt{params}。

V6 落地实现：

\begin{itemize}
\tightlist
\item
  后端新增 \texttt{queryController.executeSelectQuery}，挂载在
  \texttt{POST\ /api/v1/databases/:databaseId/query/select}。
\item
  后端按登录用户解析 \texttt{databaseId}，只允许访问当前用户拥有的物理
  schema。
\item
  查询 AST 经过体积、来源、字段、操作符、聚合函数、分页、JOIN
  数量和子查询深度校验。
\item
  \texttt{rightField} 用于 JOIN ON
  或字段间比较，普通值仍全部进入参数绑定。
\item
  顶层查询默认分页并返回
  \texttt{hasMore}，前端通过加载下一页继续追加结果。
\item
  前端 \texttt{DatabaseWorkbench.vue}
  新增``查询构造''模式，普通模式覆盖主表字段、JOIN、筛选、分组、HAVING、排序、分页；高级
  AST 面板覆盖完整子查询节点。
\end{itemize}

\hypertarget{ux5b50ux67e5ux8be2ux652fux6301ux8303ux56f4}{%
\subsubsection{10.5
子查询支持范围}\label{ux5b50ux67e5ux8be2ux652fux6301ux8303ux56f4}}

第一阶段支持：

\begin{itemize}
\tightlist
\item
  派生表：\texttt{FROM\ (SELECT\ ...)\ alias}
\item
  \texttt{IN\ (SELECT\ ...)}
\item
  \texttt{NOT\ IN\ (SELECT\ ...)}
\item
  \texttt{EXISTS\ (SELECT\ ...)}
\item
  \texttt{NOT\ EXISTS\ (SELECT\ ...)}
\item
  标量子查询：\texttt{field\ \textgreater{}\ (SELECT\ AVG(...))}
\end{itemize}

第一阶段不支持：

\begin{itemize}
\tightlist
\item
  CTE：\texttt{WITH\ ...}
\item
  窗口函数。
\item
  递归查询。
\item
  UNION。
\item
  用户输入原始 SQL 片段。
\end{itemize}

这些能力后续可以作为高级查询功能扩展。

\hypertarget{ux4e8bux52a1ux8be6ux7ec6ux8bbeux8ba1}{%
\subsection{11.
事务详细设计}\label{ux4e8bux52a1ux8be6ux7ec6ux8bbeux8ba1}}

\hypertarget{transactionservice}{%
\subsubsection{11.1 transactionService}\label{transactionservice}}

事务会话绑定后端 MySQL 连接。

核心状态：

\begin{Shaded}
\begin{Highlighting}[]
\KeywordTok{type}\NormalTok{ TransactionSession }\OperatorTok{=}\NormalTok{ \{}
\NormalTok{  id}\OperatorTok{:} \DataTypeTok{string}\OperatorTok{;}
\NormalTok{  userID}\OperatorTok{:} \DataTypeTok{number}\OperatorTok{;}
\NormalTok{  databaseId}\OperatorTok{:} \DataTypeTok{number}\OperatorTok{;}
\NormalTok{  schemaName}\OperatorTok{:} \DataTypeTok{string}\OperatorTok{;}
\NormalTok{  connectionId}\OperatorTok{:} \DataTypeTok{number}\OperatorTok{;}
\NormalTok{  status}\OperatorTok{:} \StringTok{\textquotesingle{}active\textquotesingle{}} \OperatorTok{|} \StringTok{\textquotesingle{}committed\textquotesingle{}} \OperatorTok{|} \StringTok{\textquotesingle{}rolled\_back\textquotesingle{}} \OperatorTok{|} \StringTok{\textquotesingle{}expired\textquotesingle{}}\OperatorTok{;}
\NormalTok{  isolationLevel}\OperatorTok{:} \DataTypeTok{string}\OperatorTok{;}
\NormalTok{  expiresAt}\OperatorTok{:}\NormalTok{ Date}\OperatorTok{;}
\NormalTok{\}}\OperatorTok{;}
\end{Highlighting}
\end{Shaded}

\hypertarget{ux5f00ux542fux4e8bux52a1}{%
\subsubsection{11.2 开启事务}\label{ux5f00ux542fux4e8bux52a1}}

\begin{Shaded}
\begin{Highlighting}[]
\NormalTok{1. resolveUserDatabase。}
\NormalTok{2. 从连接池获取连接。}
\NormalTok{3. USE \textasciigrave{}schemaName\textasciigrave{}。}
\NormalTok{4. SET TRANSACTION ISOLATION LEVEL ...}
\NormalTok{5. BEGIN。}
\NormalTok{6. 写入 transaction\_sessions。}
\NormalTok{7. 返回 transactionId。}
\end{Highlighting}
\end{Shaded}

\hypertarget{ux4e8bux52a1ux5185ux64cdux4f5c}{%
\subsubsection{11.3 事务内操作}\label{ux4e8bux52a1ux5185ux64cdux4f5c}}

写操作如果携带 \texttt{transactionId}：

\begin{Shaded}
\begin{Highlighting}[]
\NormalTok{1. 校验事务存在。}
\NormalTok{2. 校验事务属于当前用户。}
\NormalTok{3. 校验事务属于当前 databaseId。}
\NormalTok{4. 校验事务未过期。}
\NormalTok{5. 使用事务绑定连接执行 SQL。}
\end{Highlighting}
\end{Shaded}

\hypertarget{ux8d85ux65f6ux56deux6eda}{%
\subsubsection{11.4 超时回滚}\label{ux8d85ux65f6ux56deux6eda}}

事务超时后：

\begin{Shaded}
\begin{Highlighting}[]
\NormalTok{1. 自动 ROLLBACK。}
\NormalTok{2. 释放连接。}
\NormalTok{3. transaction\_sessions 状态改为 expired。}
\NormalTok{4. 写入审计日志。}
\end{Highlighting}
\end{Shaded}

\hypertarget{ux7edfux8ba1ux8be6ux7ec6ux8bbeux8ba1}{%
\subsection{12.
统计详细设计}\label{ux7edfux8ba1ux8be6ux7ec6ux8bbeux8ba1}}

\hypertarget{statsservice.useroverview}{%
\subsubsection{12.1
statsService.userOverview}\label{statsservice.useroverview}}

数据来源：

\begin{itemize}
\tightlist
\item
  \texttt{user\_databases}。
\item
  \texttt{information\_schema.tables}。
\item
  \texttt{audit\_logs}。
\item
  \texttt{user\_assets}。
\end{itemize}

返回：

\begin{itemize}
\tightlist
\item
  数据库数量。
\item
  表数量。
\item
  视图数量。
\item
  索引数量。
\item
  估算行数。
\item
  操作趋势。
\item
  最近操作。
\end{itemize}

\hypertarget{statsservice.databaseoverview}{%
\subsubsection{12.2
statsService.databaseOverview}\label{statsservice.databaseoverview}}

流程：

\begin{Shaded}
\begin{Highlighting}[]
\NormalTok{1. resolveUserDatabase。}
\NormalTok{2. 查询 information\_schema.tables。}
\NormalTok{3. 聚合表数量、视图数量、数据大小、索引大小。}
\NormalTok{4. 统计表大小 Top N 和行数区间分布。}
\NormalTok{5. 查询 audit\_logs 得到最近操作。}
\end{Highlighting}
\end{Shaded}

\hypertarget{statsservice.tableoverview}{%
\subsubsection{12.3
statsService.tableOverview}\label{statsservice.tableoverview}}

表级统计涉及用户数据表，必须限制字段和执行时间。

支持：

\begin{itemize}
\tightlist
\item
  行数估算。
\item
  字段数量。
\item
  主键信息。
\item
  空值比例。
\item
  数值字段 min、max、avg。
\item
  分类字段 Top N。
\end{itemize}

统计查询必须使用字段白名单，不允许前端传入聚合 SQL。

\hypertarget{ux5ba1ux8ba1ux8be6ux7ec6ux8bbeux8ba1}{%
\subsection{13.
审计详细设计}\label{ux5ba1ux8ba1ux8be6ux7ec6ux8bbeux8ba1}}

\hypertarget{auditservice}{%
\subsubsection{13.1 auditService}\label{auditservice}}

所有关键操作都调用：

\begin{Shaded}
\begin{Highlighting}[]
\NormalTok{auditService}\OperatorTok{.}\FunctionTok{record}\NormalTok{(\{}
\NormalTok{  traceId}\OperatorTok{,}
\NormalTok{  userID}\OperatorTok{,}
\NormalTok{  databaseId}\OperatorTok{,}
\NormalTok{  tableName}\OperatorTok{,}
\NormalTok{  actionType}\OperatorTok{,}
\NormalTok{  targetType}\OperatorTok{,}
\NormalTok{  summary}\OperatorTok{,}
\NormalTok{  detail}\OperatorTok{,}
\NormalTok{  success}\OperatorTok{,}
\NormalTok{  errorMessage}\OperatorTok{,}
\NormalTok{  durationMs}\OperatorTok{,}
\NormalTok{\})}\OperatorTok{;}
\end{Highlighting}
\end{Shaded}

\hypertarget{ux5ba1ux8ba1ux5185ux5bb9ux8131ux654f}{%
\subsubsection{13.2
审计内容脱敏}\label{ux5ba1ux8ba1ux5185ux5bb9ux8131ux654f}}

禁止写入：

\begin{itemize}
\tightlist
\item
  密码。
\item
  密码哈希。
\item
  Access Token。
\item
  Refresh Token。
\item
  Refresh Token 哈希。
\item
  文件绝对路径。
\item
  真实物理库名。
\end{itemize}

允许写入：

\begin{itemize}
\tightlist
\item
  用户可见数据库名。
\item
  表名。
\item
  字段名。
\item
  影响行数。
\item
  错误码。
\item
  traceId。
\end{itemize}

\hypertarget{ux524dux7aefux8be6ux7ec6ux8bbeux8ba1}{%
\subsection{14.
前端详细设计}\label{ux524dux7aefux8be6ux7ec6ux8bbeux8ba1}}

\hypertarget{ux63a8ux8350ux76eeux5f55ux7ed3ux6784-1}{%
\subsubsection{14.1
推荐目录结构}\label{ux63a8ux8350ux76eeux5f55ux7ed3ux6784-1}}

\begin{Shaded}
\begin{Highlighting}[]
\NormalTok{frontend/src/}
\NormalTok{  assets/}
\NormalTok{    theme.css}
\NormalTok{    glass.css}
\NormalTok{    global.css}

\NormalTok{  router/}
\NormalTok{    index.ts}

\NormalTok{  utils/}
\NormalTok{    request.ts}
\NormalTok{    errorMessage.ts}

\NormalTok{  api/}
\NormalTok{    authApi.ts}
\NormalTok{    workbenchApi.ts}
\NormalTok{    databaseApi.ts}
\NormalTok{    tableApi.ts}
\NormalTok{    rowApi.ts}
\NormalTok{    queryApi.ts}
\NormalTok{    statsApi.ts}
\NormalTok{    assetApi.ts}

\NormalTok{  stores/}
\NormalTok{    authStore.ts}
\NormalTok{    preferenceStore.ts}
\NormalTok{    databaseStore.ts}
\NormalTok{    workbenchStore.ts}

\NormalTok{  composables/}
\NormalTok{    useAuth.ts}
\NormalTok{    useDatabaseTree.ts}
\NormalTok{    useDataGrid.ts}
\NormalTok{    useQueryBuilder.ts}
\NormalTok{    useDashboardLayout.ts}

\NormalTok{  layouts/}
\NormalTok{    WorkbenchLayout.vue}

\NormalTok{  views/}
\NormalTok{    Login.vue}
\NormalTok{    Register.vue}
\NormalTok{    Home.vue}
\NormalTok{    DatabaseWorkbench.vue}
\NormalTok{    DatabaseDashboard.vue}
\NormalTok{    TableDesigner.vue}
\NormalTok{    TableData.vue}
\NormalTok{    QueryBuilder.vue}

\NormalTok{  components/}
\NormalTok{    ParticleBackground.vue}
\NormalTok{    AppHeader.vue}
\NormalTok{    UserAvatar.vue}
\NormalTok{    DbTree.vue}
\NormalTok{    DataGrid.vue}
\NormalTok{    DataGridToolbar.vue}
\NormalTok{    ColumnDesigner.vue}
\NormalTok{    QueryBuilderPanel.vue}
\NormalTok{    DashboardWidget.vue}
\NormalTok{    ChartCard.vue}
\NormalTok{    ConfirmDialog.vue}
\end{Highlighting}
\end{Shaded}

如果暂时不引入 Pinia，可以先用 composable + local
state；当页面状态变复杂时再引入 Pinia。正式设计推荐使用 Pinia
管理登录、偏好、数据库树和工作台布局。

\hypertarget{authstore}{%
\subsubsection{14.2 authStore}\label{authstore}}

状态：

\begin{Shaded}
\begin{Highlighting}[]
\KeywordTok{type}\NormalTok{ AuthState }\OperatorTok{=}\NormalTok{ \{}
\NormalTok{  accessToken}\OperatorTok{:} \DataTypeTok{string} \OperatorTok{|} \DataTypeTok{null}\OperatorTok{;}
\NormalTok{  user}\OperatorTok{:}\NormalTok{ \{}
\NormalTok{    userID}\OperatorTok{:} \DataTypeTok{number}\OperatorTok{;}
\NormalTok{    accountNo}\OperatorTok{:} \DataTypeTok{string}\OperatorTok{;}
\NormalTok{    userName}\OperatorTok{:} \DataTypeTok{string}\OperatorTok{;}
\NormalTok{    avatarUrl}\OperatorTok{?:} \DataTypeTok{string}\OperatorTok{;}
\NormalTok{  \} }\OperatorTok{|}\NormalTok{ null}\OperatorTok{;}
\NormalTok{\}}\OperatorTok{;}
\end{Highlighting}
\end{Shaded}

行为：

\begin{itemize}
\tightlist
\item
  login。
\item
  refreshAccessToken。
\item
  logout。
\item
  fetchMe。
\end{itemize}

\hypertarget{request.ts}{%
\subsubsection{14.3 request.ts}\label{request.ts}}

职责：

\begin{itemize}
\tightlist
\item
  自动注入 Access Token。
\item
  401 时尝试调用 \texttt{/auth/refresh}。
\item
  refresh 成功后重放原请求。
\item
  refresh 失败后清理登录状态并跳转登录页。
\item
  展示后端精准错误 message。
\end{itemize}

需要避免刷新风暴：

\begin{Shaded}
\begin{Highlighting}[]
\NormalTok{多个请求同时 401 时，只允许一个 refresh 请求进行；}
\NormalTok{其他请求等待 refresh 完成后重试。}
\end{Highlighting}
\end{Shaded}

\hypertarget{workbenchlayout}{%
\subsubsection{14.4 WorkbenchLayout}\label{workbenchlayout}}

布局：

\begin{Shaded}
\begin{Highlighting}[]
\NormalTok{顶部 AppHeader}
\NormalTok{左侧 DbTree}
\NormalTok{中间 RouterView / 动态工作区}
\NormalTok{右侧可选信息栏}
\NormalTok{背景 ParticleBackground}
\end{Highlighting}
\end{Shaded}

职责：

\begin{itemize}
\tightlist
\item
  加载用户偏好。
\item
  应用主题变量。
\item
  加载数据库树。
\item
  提供全局刷新入口。
\item
  提供退出登录入口。
\end{itemize}

\hypertarget{dbtree}{%
\subsubsection{14.5 DbTree}\label{dbtree}}

节点类型：

\begin{Shaded}
\begin{Highlighting}[]
\KeywordTok{type}\NormalTok{ DbTreeNode }\OperatorTok{=}
  \OperatorTok{|}\NormalTok{ \{ type}\OperatorTok{:} \StringTok{\textquotesingle{}database\textquotesingle{}}\OperatorTok{;}\NormalTok{ databaseId}\OperatorTok{:} \DataTypeTok{number}\OperatorTok{;}\NormalTok{ displayName}\OperatorTok{:} \DataTypeTok{string}\NormalTok{ \}}
  \OperatorTok{|}\NormalTok{ \{ type}\OperatorTok{:} \StringTok{\textquotesingle{}table\textquotesingle{}}\OperatorTok{;}\NormalTok{ databaseId}\OperatorTok{:} \DataTypeTok{number}\OperatorTok{;}\NormalTok{ tableName}\OperatorTok{:} \DataTypeTok{string}\NormalTok{ \}}
  \OperatorTok{|}\NormalTok{ \{ type}\OperatorTok{:} \StringTok{\textquotesingle{}view\textquotesingle{}}\OperatorTok{;}\NormalTok{ databaseId}\OperatorTok{:} \DataTypeTok{number}\OperatorTok{;}\NormalTok{ viewName}\OperatorTok{:} \DataTypeTok{string}\NormalTok{ \}}\OperatorTok{;}
\end{Highlighting}
\end{Shaded}

交互：

\begin{itemize}
\tightlist
\item
  点击数据库打开数据库看板。
\item
  点击表打开表数据页。
\item
  右键数据库：创建表、刷新、删除数据库。
\item
  右键表：查看结构、打开数据、删除表。
\end{itemize}

\hypertarget{datagrid}{%
\subsubsection{14.6 DataGrid}\label{datagrid}}

职责：

\begin{itemize}
\tightlist
\item
  显示 \texttt{API-ROW-001} 返回的 columns 和 items。
\item
  显示预览接口返回的 columns、rows、total 和 hasMore。
\item
  支持分页预览和字段筛选。
\item
  支持行内编辑草稿态。
\item
  支持单行更新、取消。
\item
  支持新增行。
\item
  支持基于主键删除单行。
\item
  支持排序、列显示控制、JSON/CSV
  批量新增、基于主键列表的批量修改和批量删除。
\end{itemize}

行状态：

\begin{Shaded}
\begin{Highlighting}[]
\KeywordTok{type}\NormalTok{ DataGridRowState }\OperatorTok{=}\NormalTok{ \{}
\NormalTok{  rowKey}\OperatorTok{:} \DataTypeTok{string}\OperatorTok{;}
\NormalTok{  original}\OperatorTok{:} \BuiltInTok{Record}\OperatorTok{\textless{}}\DataTypeTok{string}\OperatorTok{,} \DataTypeTok{unknown}\OperatorTok{\textgreater{};}
\NormalTok{  draft}\OperatorTok{:} \BuiltInTok{Record}\OperatorTok{\textless{}}\DataTypeTok{string}\OperatorTok{,} \DataTypeTok{unknown}\OperatorTok{\textgreater{};}
\NormalTok{  changedFields}\OperatorTok{:} \BuiltInTok{Record}\OperatorTok{\textless{}}\DataTypeTok{string}\OperatorTok{,} \DataTypeTok{unknown}\OperatorTok{\textgreater{};}
\NormalTok{  status}\OperatorTok{:} \StringTok{\textquotesingle{}clean\textquotesingle{}} \OperatorTok{|} \StringTok{\textquotesingle{}dirty\textquotesingle{}} \OperatorTok{|} \StringTok{\textquotesingle{}saving\textquotesingle{}} \OperatorTok{|} \StringTok{\textquotesingle{}error\textquotesingle{}}\OperatorTok{;}
\NormalTok{\}}\OperatorTok{;}
\end{Highlighting}
\end{Shaded}

主键生成 rowKey：

\begin{Shaded}
\begin{Highlighting}[]
\NormalTok{单主键：String(row[pk])}
\NormalTok{复合主键：JSON.stringify(pkValues)}
\NormalTok{无主键：只读模式}
\end{Highlighting}
\end{Shaded}

\hypertarget{querybuilderpanel}{%
\subsubsection{14.7 QueryBuilderPanel}\label{querybuilderpanel}}

职责：

\begin{itemize}
\tightlist
\item
  通过 GUI 构建查询 AST。
\item
  支持主表、JOIN、字段、条件、分组、聚合、排序、分页。
\item
  支持子查询节点编辑。
\item
  在前端做基础结构校验。
\item
  提交给后端做最终校验和执行。
\end{itemize}

前端只负责构建 AST，不负责拼 SQL。

\hypertarget{dashboardwidget}{%
\subsubsection{14.8 DashboardWidget}\label{dashboardwidget}}

工作台卡片类型：

\begin{Shaded}
\begin{Highlighting}[]
\NormalTok{stat{-}number}
\NormalTok{line{-}chart}
\NormalTok{bar{-}chart}
\NormalTok{pie{-}chart}
\NormalTok{recent{-}operations}
\NormalTok{database{-}shortcut}
\NormalTok{quick{-}action}
\end{Highlighting}
\end{Shaded}

布局状态：

\begin{Shaded}
\begin{Highlighting}[]
\KeywordTok{type}\NormalTok{ DashboardWidgetState }\OperatorTok{=}\NormalTok{ \{}
\NormalTok{  id}\OperatorTok{:} \DataTypeTok{string}\OperatorTok{;}
\NormalTok{  widgetType}\OperatorTok{:} \DataTypeTok{string}\OperatorTok{;}
\NormalTok{  title}\OperatorTok{:} \DataTypeTok{string}\OperatorTok{;}
\NormalTok{  x}\OperatorTok{:} \DataTypeTok{number}\OperatorTok{;}
\NormalTok{  y}\OperatorTok{:} \DataTypeTok{number}\OperatorTok{;}
\NormalTok{  width}\OperatorTok{:} \DataTypeTok{number}\OperatorTok{;}
\NormalTok{  height}\OperatorTok{:} \DataTypeTok{number}\OperatorTok{;}
\NormalTok{  config}\OperatorTok{:} \BuiltInTok{Record}\OperatorTok{\textless{}}\DataTypeTok{string}\OperatorTok{,} \DataTypeTok{unknown}\OperatorTok{\textgreater{};}
\NormalTok{  visible}\OperatorTok{:} \DataTypeTok{boolean}\OperatorTok{;}
\NormalTok{\}}\OperatorTok{;}
\end{Highlighting}
\end{Shaded}

第一阶段可以先支持显示、隐藏、排序；拖拽布局作为后续增强。

\hypertarget{ux524dux540eux7aefux4ea4ux4e92ux6d41ux7a0b}{%
\subsection{15.
前后端交互流程}\label{ux524dux540eux7aefux4ea4ux4e92ux6d41ux7a0b}}

\hypertarget{ux767bux5f55ux540eux8fdbux5165ux5de5ux4f5cux53f0}{%
\subsubsection{15.1
登录后进入工作台}\label{ux767bux5f55ux540eux8fdbux5165ux5de5ux4f5cux53f0}}

\begin{Shaded}
\begin{Highlighting}[]
\NormalTok{Login.vue}
\NormalTok{  {-}\textgreater{} authApi.login}
\NormalTok{  {-}\textgreater{} 保存 accessToken}
\NormalTok{  {-}\textgreater{} router.push(\textquotesingle{}/home\textquotesingle{})}

\NormalTok{Home.vue / WorkbenchLayout}
\NormalTok{  {-}\textgreater{} authApi.me}
\NormalTok{  {-}\textgreater{} preferenceApi.get}
\NormalTok{  {-}\textgreater{} workbenchApi.overview}
\NormalTok{  {-}\textgreater{} databaseApi.list}
\NormalTok{  {-}\textgreater{} 渲染工作台}
\end{Highlighting}
\end{Shaded}

\hypertarget{ux521bux5efaux6570ux636eux5e93}{%
\subsubsection{15.2 创建数据库}\label{ux521bux5efaux6570ux636eux5e93}}

\begin{Shaded}
\begin{Highlighting}[]
\NormalTok{用户点击创建数据库}
\NormalTok{  {-}\textgreater{} 前端校验名称格式}
\NormalTok{  {-}\textgreater{} POST /databases}
\NormalTok{  {-}\textgreater{} 后端校验 token}
\NormalTok{  {-}\textgreater{} 后端检查容量限制}
\NormalTok{  {-}\textgreater{} 后端 CREATE DATABASE}
\NormalTok{  {-}\textgreater{} 后端写 user\_databases}
\NormalTok{  {-}\textgreater{} 后端写 audit\_logs}
\NormalTok{  {-}\textgreater{} 前端刷新 DbTree}
\end{Highlighting}
\end{Shaded}

\hypertarget{ux884cux5185ux7f16ux8f91ux63d0ux4ea4}{%
\subsubsection{15.3
行内编辑提交}\label{ux884cux5185ux7f16ux8f91ux63d0ux4ea4}}

\begin{Shaded}
\begin{Highlighting}[]
\NormalTok{用户编辑单元格}
\NormalTok{  {-}\textgreater{} DataGrid 标记 dirty}
\NormalTok{  {-}\textgreater{} 用户点击行级更新}
\NormalTok{  {-}\textgreater{} PATCH /rows \{ primaryKey, set \}}
\NormalTok{  {-}\textgreater{} 后端校验归属、表、字段、主键}
\NormalTok{  {-}\textgreater{} 后端 UPDATE}
\NormalTok{  {-}\textgreater{} 后端写 audit\_logs}
\NormalTok{  {-}\textgreater{} 前端更新 original}
\end{Highlighting}
\end{Shaded}

\hypertarget{ux7ed3ux6784ux5316ux67e5ux8be2}{%
\subsubsection{15.4 结构化查询}\label{ux7ed3ux6784ux5316ux67e5ux8be2}}

\begin{Shaded}
\begin{Highlighting}[]
\NormalTok{用户在 QueryBuilder 中配置查询}
\NormalTok{  {-}\textgreater{} 前端生成 AST}
\NormalTok{  {-}\textgreater{} 前端基础校验}
\NormalTok{  {-}\textgreater{} POST /query/select}
\NormalTok{  {-}\textgreater{} 后端 AST 校验}
\NormalTok{  {-}\textgreater{} 后端作用域解析}
\NormalTok{  {-}\textgreater{} 后端生成 SELECT + params}
\NormalTok{  {-}\textgreater{} 后端执行}
\NormalTok{  {-}\textgreater{} 返回 columns、items、pagination}
\end{Highlighting}
\end{Shaded}

\hypertarget{ux5f00ux53d1ux843dux5730ux987aux5e8f}{%
\subsection{16.
开发落地顺序}\label{ux5f00ux53d1ux843dux5730ux987aux5e8f}}

\hypertarget{phase-1ux8ba4ux8bc1ux4e0eux5de5ux4f5cux53f0ux5730ux57fa}{%
\subsubsection{Phase
1：认证与工作台地基}\label{phase-1ux8ba4ux8bc1ux4e0eux5de5ux4f5cux53f0ux5730ux57fa}}

\begin{itemize}
\tightlist
\item
  \texttt{env.ts} 完整化。
\item
  \texttt{traceMiddleware}。
\item
  \texttt{errorMiddleware}。
\item
  双 Token 设计中的 Access Token 先落地，Refresh Token 可紧随其后。
\item
  \texttt{GET\ /auth/me}。
\item
  \texttt{Home.vue} / \texttt{WorkbenchLayout} 骨架。
\end{itemize}

\hypertarget{phase-2ux7528ux6237ux504fux597dux548cux8d44ux6e90}{%
\subsubsection{Phase
2：用户偏好和资源}\label{phase-2ux7528ux6237ux504fux597dux548cux8d44ux6e90}}

\begin{itemize}
\tightlist
\item
  \texttt{user\_profiles}。
\item
  \texttt{user\_preferences}。
\item
  头像和背景资源元信息。
\item
  主题色和背景配置。
\end{itemize}

\hypertarget{phase-3ux6570ux636eux5e93ux7ba1ux7406}{%
\subsubsection{Phase
3：数据库管理}\label{phase-3ux6570ux636eux5e93ux7ba1ux7406}}

\begin{itemize}
\tightlist
\item
  \texttt{user\_databases}。
\item
  数据库列表。
\item
  创建数据库。
\item
  删除数据库。
\item
  容量限制。
\item
  审计日志。
\end{itemize}

\hypertarget{phase-4ux8868ux7ed3ux6784ux7ba1ux7406}{%
\subsubsection{Phase
4：表结构管理}\label{phase-4ux8868ux7ed3ux6784ux7ba1ux7406}}

\begin{itemize}
\tightlist
\item
  表列表。
\item
  创建表。
\item
  表结构查看。
\item
  删除表。
\item
  字段、索引和约束管理。
\item
  视图逐步扩展。
\end{itemize}

当前实现范围：

\begin{itemize}
\tightlist
\item
  已完成数据库对象列表、结构化创建表、表结构读取、删除表和在线结构管理。
\item
  创建表阶段支持字段、可选主键、唯一约束、外键、有限取值检查约束和普通/唯一索引的服务端校验与
  DDL 生成。
\item
  在线结构管理支持字段新增、字段修改、字段删除、索引新增/删除、约束新增/删除，其中有限取值通过
  CHECK 约束表达。
\item
  视图管理仍随 GUI 查询构造器后续落地。
\end{itemize}

\hypertarget{phase-5ux8868ux6570ux636e-crud}{%
\subsubsection{Phase 5：表数据
CRUD}\label{phase-5ux8868ux6570ux636e-crud}}

\begin{itemize}
\tightlist
\item
  分页预览和字段筛选。
\item
  排序和列显示控制。
\item
  单行新增。
\item
  批量新增。
\item
  行内编辑草稿态。
\item
  基于主键的行级更新。
\item
  基于主键列表的批量修改。
\item
  基于主键的单行删除。
\item
  基于主键列表的批量删除。
\item
  无主键表浏览、新增和只读保护。
\item
  基于主键列表的批量写入限制、二次确认和事务原子性。
\end{itemize}

\hypertarget{phase-6gui-ux67e5ux8be2ux6784ux9020ux5668}{%
\subsubsection{Phase 6：GUI
查询构造器}\label{phase-6gui-ux67e5ux8be2ux6784ux9020ux5668}}

\begin{itemize}
\tightlist
\item
  基础 SELECT：已落地。
\item
  JOIN：已落地普通表单和后端校验。
\item
  聚合：已落地
  \texttt{COUNT}、\texttt{SUM}、\texttt{AVG}、\texttt{MIN}、\texttt{MAX}
  白名单。
\item
  分组：已落地。
\item
  HAVING：已落地普通表单和后端校验。
\item
  子查询：已落地后端 AST，前端通过高级 AST
  面板表达派生表、\texttt{IN}、\texttt{EXISTS} 和标量子查询。
\item
  查询复杂度限制：已落地 AST 体积、JOIN
  数、子查询深度、分页大小和执行时间提示。
\end{itemize}

\hypertarget{phase-7ux7edfux8ba1ux548c-diy-ux5de5ux4f5cux53f0}{%
\subsubsection{Phase 7：统计和 DIY
工作台}\label{phase-7ux7edfux8ba1ux548c-diy-ux5de5ux4f5cux53f0}}

V7 首轮目标是``可用的统计
Dashboard''，而不是一次性做成无限自定义门户。实现时必须保证现有
Home、数据库工作台、表数据预览和 GUI 查询能力不回退。

小阶段：

\begin{longtable}[]{@{}
  >{\raggedright\arraybackslash}p{(\columnwidth - 6\tabcolsep) * \real{0.25}}
  >{\raggedright\arraybackslash}p{(\columnwidth - 6\tabcolsep) * \real{0.25}}
  >{\raggedright\arraybackslash}p{(\columnwidth - 6\tabcolsep) * \real{0.25}}
  >{\raggedright\arraybackslash}p{(\columnwidth - 6\tabcolsep) * \real{0.25}}@{}}
\toprule
小阶段 & 后端交付 & 前端交付 & 验收重点 \\
\midrule
\endhead
V7.1 & \texttt{statsService}、\texttt{auditService.list}、统计与审计路由
& 统计接口工具函数、审计列表数据模型 &
接口鉴权、归属校验、分页、容量限制 \\
V7.2 & 用户级统计聚合 & 首页统计卡片、操作趋势、数据库分布、最近操作 &
首页只放总览、入口和统计，不放表内部操作 \\
V7.3 & 数据库级和表级统计聚合 & 工作台统计入口、表结构窗口的统计区域 &
大表不阻塞，空表和无权限错误可读 \\
V7.4 & 偏好接口支持 \texttt{dashboardLayout} &
卡片显示、隐藏、排序和保存 & 刷新后布局保持，主题色全局生效 \\
V7.5 & 统一错误、审计、性能收口 & 响应式、加载态、空状态、错误态统一 &
构建通过，横向滚动尽量避免 \\
\bottomrule
\end{longtable}

首轮纳入：

\begin{itemize}
\tightlist
\item
  用户级统计：数据库数量、表数量、视图数量、索引数量、估算行数、容量占用、操作趋势、最近操作。
\item
  数据库级统计：表/视图/索引数量、数据大小、索引大小、表大小 Top
  N、表行数分布。
\item
  表级统计：字段数量、主键、外键、索引、估算行数、空值比例、数值字段聚合、分类字段
  Top N。
\item
  审计日志列表：分页、时间范围、对象类型、操作类型、数据库过滤、traceId
  检索。
\item
  Dashboard DIY：卡片显示、隐藏、排序和卡片级配置持久化。
\end{itemize}

首轮不纳入：

\begin{itemize}
\tightlist
\item
  自由拖拽网格。
\item
  多套布局方案切换。
\item
  用户自定义脚本卡片。
\item
  后台定时统计任务。
\item
  对超大表执行无上限 \texttt{COUNT(*)} 或全表字段扫描。
\end{itemize}

统计实现原则：

\begin{itemize}
\tightlist
\item
  用户级和数据库级优先使用 \texttt{information\_schema} 估算值。
\item
  表级聚合必须限制字段数量、Top N 数量、采样行数或执行时间。
\item
  所有字段名、表名、排序字段必须来自后端白名单，不接受前端传 SQL 片段。
\item
  统计访问写入审计日志，但审计详情不保存完整查询参数中的敏感值。
\end{itemize}

\hypertarget{ux5b9eux73b0ux68c0ux67e5ux6e05ux5355}{%
\subsection{17.
实现检查清单}\label{ux5b9eux73b0ux68c0ux67e5ux6e05ux5355}}

后端实现一个接口时必须检查：

\begin{itemize}
\tightlist
\item
  是否经过 authMiddleware。
\item
  是否生成 traceId。
\item
  是否校验数据库归属。
\item
  是否校验标识符。
\item
  是否使用 SQL builder。
\item
  是否参数化字段值。
\item
  是否执行容量限制。
\item
  是否记录审计日志。
\item
  是否返回统一响应。
\item
  是否映射 MySQL 错误。
\end{itemize}

前端实现一个页面时必须检查：

\begin{itemize}
\tightlist
\item
  是否使用 request 统一请求。
\item
  是否处理 loading。
\item
  是否展示后端精准错误。
\item
  是否处理 401 自动刷新。
\item
  是否避免直接拼 SQL。
\item
  是否把待提交修改和已保存数据区分。
\item
  是否给高风险操作二次确认。
\end{itemize}

\hypertarget{ux975eux529fux80fdux6027ux8bbeux8ba1}{%
\subsection{18.
非功能性设计}\label{ux975eux529fux80fdux6027ux8bbeux8ba1}}

\hypertarget{ux6027ux80fdux8bbeux8ba1}{%
\subsubsection{18.1 性能设计}\label{ux6027ux80fdux8bbeux8ba1}}

\begin{itemize}
\tightlist
\item
  所有列表接口必须分页，默认 \texttt{pageSize\ =\ 20}，最大
  \texttt{pageSize\ =\ 100}。
\item
  查询构造器最大返回 \texttt{1000} 行，超过必须分页或缩小条件。
\item
  查询最大执行时间默认 \texttt{10} 秒，超时后返回容量或超时错误。
\item
  统计接口优先使用 \texttt{information\_schema}
  的估算数据，避免对大表频繁执行 \texttt{COUNT(*)}。
\item
  表格数据页只加载当前页，不一次性拉取整表。
\end{itemize}

\hypertarget{ux53efux7528ux6027ux8bbeux8ba1}{%
\subsubsection{18.2 可用性设计}\label{ux53efux7528ux6027ux8bbeux8ba1}}

\begin{itemize}
\tightlist
\item
  MySQL 不可用时，数据库管理能力降级为明确错误提示，前端不能白屏。
\item
  Access Token 过期时，前端先尝试刷新 Token；刷新失败后才跳转登录页。
\item
  高风险操作必须给出影响说明和确认方式。
\item
  表格行内编辑必须支持取消，避免误提交。
\end{itemize}

\hypertarget{ux53efux7ef4ux62a4ux6027ux8bbeux8ba1}{%
\subsubsection{18.3
可维护性设计}\label{ux53efux7ef4ux62a4ux6027ux8bbeux8ba1}}

\begin{itemize}
\tightlist
\item
  controller 只处理请求和响应，不拼接 SQL。
\item
  SQL 构造集中在 \texttt{sql/} 目录。
\item
  错误响应集中在 \texttt{errorMiddleware}。
\item
  容量限制集中在 \texttt{limits.ts} 和 \texttt{limitService}。
\item
  前端 API 调用集中在 \texttt{api/} 目录。
\end{itemize}

\hypertarget{ux53efux89c2ux6d4bux6027ux8bbeux8ba1}{%
\subsubsection{18.4
可观测性设计}\label{ux53efux89c2ux6d4bux6027ux8bbeux8ba1}}

\begin{itemize}
\tightlist
\item
  每个请求都有 \texttt{traceId}。
\item
  关键操作写入 \texttt{audit\_logs}。
\item
  后端日志、接口响应、审计日志都包含同一个 \texttt{traceId}。
\item
  MySQL 错误必须映射为可读错误码。
\end{itemize}

\hypertarget{ux53efux6269ux5c55ux6027ux8bbeux8ba1}{%
\subsubsection{18.5
可扩展性设计}\label{ux53efux6269ux5c55ux6027ux8bbeux8ba1}}

\begin{itemize}
\tightlist
\item
  用户物理库通过 \texttt{user\_databases} 映射，不把真实库名暴露给前端。
\item
  工作台卡片使用 \texttt{widgetType\ +\ config}
  模型，后续可以扩展新卡片。
\item
  查询构造器使用 AST，后续可以扩展更多查询节点。
\end{itemize}

\hypertarget{ux9632ux5fa1ux6027ux8bbeux8ba1}{%
\subsection{19. 防御性设计}\label{ux9632ux5fa1ux6027ux8bbeux8ba1}}

\hypertarget{ux8f93ux5165ux9632ux5fa1}{%
\subsubsection{19.1 输入防御}\label{ux8f93ux5165ux9632ux5fa1}}

\begin{itemize}
\tightlist
\item
  所有请求参数必须经过 controller 或 validator 校验。
\item
  数据库名、表名、字段名、索引名、约束名、视图名必须符合标识符规则。
\item
  文件上传必须校验 MIME、扩展名、大小和文件头。
\end{itemize}

\hypertarget{sql-ux9632ux5fa1}{%
\subsubsection{19.2 SQL 防御}\label{sql-ux9632ux5fa1}}

\begin{itemize}
\tightlist
\item
  字段值只能通过参数化传入。
\item
  标识符只能来自校验后的用户输入或 MySQL 元数据。
\item
  业务代码不能手写反引号拼接 SQL。
\item
  查询构造器只生成 \texttt{SELECT}，不生成写操作 SQL。
\end{itemize}

\hypertarget{ux6743ux9650ux9632ux5fa1}{%
\subsubsection{19.3 权限防御}\label{ux6743ux9650ux9632ux5fa1}}

\begin{itemize}
\tightlist
\item
  所有数据库相关接口都先执行 \texttt{databaseResolver}。
\item
  事务、资源、偏好、审计查询都必须校验归属。
\item
  不向前端返回真实物理库名、文件绝对路径、Token 哈希。
\end{itemize}

\hypertarget{ux72b6ux6001ux9632ux5fa1}{%
\subsubsection{19.4 状态防御}\label{ux72b6ux6001ux9632ux5fa1}}

\begin{itemize}
\tightlist
\item
  Refresh Token 使用轮换机制。
\item
  事务必须有超时时间，超时自动回滚。
\item
  创建数据库时如果系统表写入失败，必须尝试补偿删除物理库。
\item
  行内编辑必须先进入草稿态，只有显式提交才写数据库。
\end{itemize}

\hypertarget{ux5bb9ux91cfux9632ux5fa1}{%
\subsubsection{19.5 容量防御}\label{ux5bb9ux91cfux9632ux5fa1}}

\begin{itemize}
\tightlist
\item
  创建类操作先检查容量，再执行写入。
\item
  查询类操作检查 AST 大小、JOIN 数、子查询深度、分页大小和执行时间。
\item
  批量修改和删除先预估影响行数，超过阈值要求确认或拒绝。
\end{itemize}

\hypertarget{ux8be6ux7ec6ux8bbeux8ba1ux8ffdux8e2a}{%
\subsection{20.
详细设计追踪}\label{ux8be6ux7ec6ux8bbeux8ba1ux8ffdux8e2a}}

为避免追踪信息在 PDF 宽表中挤压，详细设计追踪按设计项列出：

\begin{itemize}
\tightlist
\item
  双 Token：对应需求 \texttt{REQ-AUTH}；对应接口
  \texttt{API-AUTH}；对应安全 \texttt{SEC-AUTH}。
\item
  工作台偏好：对应需求
  \texttt{REQ-WORKBENCH}、\texttt{REQ-PREF}；对应接口
  \texttt{API-WORKBENCH}、\texttt{API-PREF}；对应安全
  \texttt{SEC-AUTHZ}。
\item
  用户资源：对应需求 \texttt{REQ-ASSET}；对应接口
  \texttt{API-ASSET}；对应安全 \texttt{SEC-FILE}。
\item
  数据库管理：对应需求 \texttt{REQ-DB}；对应接口
  \texttt{API-DB}；对应安全 \texttt{SEC-AUTHZ}、\texttt{SEC-SQL}。
\item
  表结构管理：对应需求 \texttt{REQ-TABLE}；对应接口
  \texttt{API-TABLE}；对应安全 \texttt{SEC-SQL}、\texttt{SEC-CONFIRM}。
\item
  行内编辑：对应需求 \texttt{REQ-ROW}；对应接口
  \texttt{API-ROW}；对应安全 \texttt{SEC-SQL}、\texttt{SEC-CONFIRM}。
\item
  查询构造器：对应需求 \texttt{REQ-QUERY}；对应接口
  \texttt{API-QUERY}；对应安全 \texttt{SEC-SQL}、\texttt{SEC-LIMIT}。
\item
  事务：对应需求 \texttt{REQ-TX}；对应接口 \texttt{API-TX}；对应安全
  \texttt{SEC-TX}。
\item
  统计：对应需求 \texttt{REQ-STATS}；对应接口
  \texttt{API-STATS}；对应安全 \texttt{SEC-AUTHZ}、\texttt{SEC-LIMIT}。
\item
  审计：对应需求 \texttt{REQ-AUDIT}；对应接口
  \texttt{API-AUDIT}；对应安全 \texttt{SEC-AUDIT}。
\item
  容量限制：对应需求 \texttt{REQ-LIMIT}；对应接口规范
  \texttt{2.4}；对应安全 \texttt{SEC-LIMIT}。
\end{itemize}

\hypertarget{ux6570ux636eux5e93ux8bbeux8ba1ux9644ux5f55}{%
\section{数据库设计附录}\label{ux6570ux636eux5e93ux8bbeux8ba1ux9644ux5f55}}

\hypertarget{ux6570ux636eux5e93ux8bbeux8ba1}{%
\section{数据库设计}\label{ux6570ux636eux5e93ux8bbeux8ba1}}

\hypertarget{ux8bbeux8ba1ux539fux5219}{%
\subsection{1. 设计原则}\label{ux8bbeux8ba1ux539fux5219}}

\hypertarget{data-001-ux7cfbux7edfux5e93ux548cux7528ux6237ux7269ux7406ux5e93ux5206ux79bb}{%
\subsubsection{DATA-001
系统库和用户物理库分离}\label{data-001-ux7cfbux7edfux5e93ux548cux7528ux6237ux7269ux7406ux5e93ux5206ux79bb}}

系统库 \texttt{DB\_App}
保存平台元数据。用户物理库保存用户自行创建的真实数据。

\hypertarget{data-002-ux7b2cux4e09ux8303ux5f0f}{%
\subsubsection{DATA-002
第三范式}\label{data-002-ux7b2cux4e09ux8303ux5f0f}}

系统库核心表满足第三范式：

\begin{itemize}
\tightlist
\item
  \texttt{users} 只保存认证信息。
\item
  \texttt{user\_profiles} 保存展示资料。
\item
  \texttt{user\_preferences} 保存偏好设置。
\item
  \texttt{user\_assets} 保存资源元信息。
\item
  \texttt{user\_databases} 保存数据库归属关系。
\item
  \texttt{refresh\_tokens} 保存会话刷新信息。
\item
  \texttt{audit\_logs} 保存操作追踪。
\end{itemize}

\hypertarget{data-003-ux7528ux6237ux8868ux7ed3ux6784ux81eaux7531}{%
\subsubsection{DATA-003
用户表结构自由}\label{data-003-ux7528ux6237ux8868ux7ed3ux6784ux81eaux7531}}

用户物理库中的表由用户通过 GUI
自行设计。系统只限制安全、容量和权限，不限制用户表的业务字段。

\hypertarget{data-004-ux771fux5b9eux7269ux7406ux5e93ux540dux540eux7aefux751fux6210}{%
\subsubsection{DATA-004
真实物理库名后端生成}\label{data-004-ux771fux5b9eux7269ux7406ux5e93ux540dux540eux7aefux751fux6210}}

前端只使用平台数据库 ID。真实 MySQL database 名只由后端生成和使用。

\hypertarget{ux7cfbux7edfux5e93ux8868ux603bux89c8}{%
\subsection{2.
系统库表总览}\label{ux7cfbux7edfux5e93ux8868ux603bux89c8}}

\begin{longtable}[]{@{}ll@{}}
\toprule
表名 & 说明 \\
\midrule
\endhead
users & 用户认证表 \\
user\_profiles & 用户资料表 \\
user\_preferences & 用户偏好表 \\
user\_assets & 用户资源元信息表 \\
refresh\_tokens & Refresh Token 表 \\
user\_databases & 用户数据库归属表 \\
transaction\_sessions & 事务会话表 \\
audit\_logs & 操作审计日志表 \\
stat\_snapshots & 统计快照表 \\
\bottomrule
\end{longtable}

\hypertarget{ux547dux540dux89c4ux5219}{%
\subsection{3. 命名规则}\label{ux547dux540dux89c4ux5219}}

数据库显示名、表名、字段名、索引名、约束名、视图名统一使用：

\begin{Shaded}
\begin{Highlighting}[]
\NormalTok{\^{}[A{-}Za{-}z][A{-}Za{-}z0{-}9\_]\{0,63\}$}
\end{Highlighting}
\end{Shaded}

规则：

\begin{itemize}
\tightlist
\item
  必须以英文字母开头。
\item
  只允许英文字母、数字、下划线。
\item
  最长 64 字符。
\item
  不允许空格、中文、横线、点号、反引号。
\end{itemize}

\hypertarget{users-ux7528ux6237ux8ba4ux8bc1ux8868}{%
\subsection{4. users
用户认证表}\label{users-ux7528ux6237ux8ba4ux8bc1ux8868}}

\begin{Shaded}
\begin{Highlighting}[]
\KeywordTok{CREATE} \KeywordTok{TABLE} \ControlFlowTok{IF} \KeywordTok{NOT} \KeywordTok{EXISTS}\NormalTok{ users (}
  \KeywordTok{id} \DataTypeTok{INT}\NormalTok{ AUTO\_INCREMENT }\KeywordTok{PRIMARY} \KeywordTok{KEY} \KeywordTok{COMMENT} \StringTok{\textquotesingle{}用户内部唯一 ID\textquotesingle{}}\NormalTok{,}
\NormalTok{  user\_name }\DataTypeTok{VARCHAR}\NormalTok{(}\DecValTok{50}\NormalTok{) }\KeywordTok{NOT} \KeywordTok{NULL} \KeywordTok{COMMENT} \StringTok{\textquotesingle{}用户名，仅供展示\textquotesingle{}}\NormalTok{,}
\NormalTok{  account\_no }\DataTypeTok{VARCHAR}\NormalTok{(}\DecValTok{10}\NormalTok{) }\KeywordTok{NOT} \KeywordTok{NULL} \KeywordTok{COMMENT} \StringTok{\textquotesingle{}登录账号，10 位随机数字\textquotesingle{}}\NormalTok{,}
\NormalTok{  password\_hash }\DataTypeTok{VARCHAR}\NormalTok{(}\DecValTok{255}\NormalTok{) }\KeywordTok{NOT} \KeywordTok{NULL} \KeywordTok{COMMENT} \StringTok{\textquotesingle{}bcrypt 密码哈希值\textquotesingle{}}\NormalTok{,}
\NormalTok{  status ENUM(}\StringTok{\textquotesingle{}active\textquotesingle{}}\NormalTok{, }\StringTok{\textquotesingle{}disabled\textquotesingle{}}\NormalTok{) }\KeywordTok{NOT} \KeywordTok{NULL} \KeywordTok{DEFAULT} \StringTok{\textquotesingle{}active\textquotesingle{}} \KeywordTok{COMMENT} \StringTok{\textquotesingle{}用户状态\textquotesingle{}}\NormalTok{,}
\NormalTok{  register\_time }\DataTypeTok{TIMESTAMP} \KeywordTok{NOT} \KeywordTok{NULL} \KeywordTok{DEFAULT} \FunctionTok{CURRENT\_TIMESTAMP} \KeywordTok{COMMENT} \StringTok{\textquotesingle{}账号注册时间\textquotesingle{}}\NormalTok{,}
\NormalTok{  last\_login\_at }\DataTypeTok{TIMESTAMP} \KeywordTok{NULL} \KeywordTok{DEFAULT} \KeywordTok{NULL} \KeywordTok{COMMENT} \StringTok{\textquotesingle{}最近登录时间\textquotesingle{}}\NormalTok{,}

  \KeywordTok{UNIQUE} \KeywordTok{KEY}\NormalTok{ uk\_users\_account\_no (account\_no),}
  \KeywordTok{KEY}\NormalTok{ idx\_users\_status (status)}
\NormalTok{) }\KeywordTok{COMMENT}\OperatorTok{=}\StringTok{\textquotesingle{}系统用户表\textquotesingle{}}\NormalTok{;}
\end{Highlighting}
\end{Shaded}

\hypertarget{user_profiles-ux7528ux6237ux8d44ux6599ux8868}{%
\subsection{5. user\_profiles
用户资料表}\label{user_profiles-ux7528ux6237ux8d44ux6599ux8868}}

\begin{Shaded}
\begin{Highlighting}[]
\KeywordTok{CREATE} \KeywordTok{TABLE} \ControlFlowTok{IF} \KeywordTok{NOT} \KeywordTok{EXISTS}\NormalTok{ user\_profiles (}
\NormalTok{  user\_id }\DataTypeTok{INT} \KeywordTok{PRIMARY} \KeywordTok{KEY} \KeywordTok{COMMENT} \StringTok{\textquotesingle{}用户 ID\textquotesingle{}}\NormalTok{,}
\NormalTok{  nickname }\DataTypeTok{VARCHAR}\NormalTok{(}\DecValTok{50}\NormalTok{) }\KeywordTok{NOT} \KeywordTok{NULL} \KeywordTok{COMMENT} \StringTok{\textquotesingle{}展示昵称\textquotesingle{}}\NormalTok{,}
\NormalTok{  avatar\_asset\_id BIGINT }\KeywordTok{NULL} \KeywordTok{COMMENT} \StringTok{\textquotesingle{}头像资源 ID\textquotesingle{}}\NormalTok{,}
\NormalTok{  created\_at }\DataTypeTok{TIMESTAMP} \KeywordTok{NOT} \KeywordTok{NULL} \KeywordTok{DEFAULT} \FunctionTok{CURRENT\_TIMESTAMP} \KeywordTok{COMMENT} \StringTok{\textquotesingle{}创建时间\textquotesingle{}}\NormalTok{,}
\NormalTok{  updated\_at }\DataTypeTok{TIMESTAMP} \KeywordTok{NOT} \KeywordTok{NULL} \KeywordTok{DEFAULT} \FunctionTok{CURRENT\_TIMESTAMP} \KeywordTok{ON} \KeywordTok{UPDATE} \FunctionTok{CURRENT\_TIMESTAMP} \KeywordTok{COMMENT} \StringTok{\textquotesingle{}更新时间\textquotesingle{}}\NormalTok{,}

  \KeywordTok{CONSTRAINT}\NormalTok{ fk\_profiles\_user}
    \KeywordTok{FOREIGN} \KeywordTok{KEY}\NormalTok{ (user\_id) }\KeywordTok{REFERENCES}\NormalTok{ users(}\KeywordTok{id}\NormalTok{)}
    \KeywordTok{ON} \KeywordTok{DELETE} \KeywordTok{CASCADE}
\NormalTok{) }\KeywordTok{COMMENT}\OperatorTok{=}\StringTok{\textquotesingle{}用户资料表\textquotesingle{}}\NormalTok{;}
\end{Highlighting}
\end{Shaded}

\hypertarget{user_preferences-ux7528ux6237ux504fux597dux8868}{%
\subsection{6. user\_preferences
用户偏好表}\label{user_preferences-ux7528ux6237ux504fux597dux8868}}

\begin{Shaded}
\begin{Highlighting}[]
\KeywordTok{CREATE} \KeywordTok{TABLE} \ControlFlowTok{IF} \KeywordTok{NOT} \KeywordTok{EXISTS}\NormalTok{ user\_preferences (}
\NormalTok{  user\_id }\DataTypeTok{INT} \KeywordTok{PRIMARY} \KeywordTok{KEY} \KeywordTok{COMMENT} \StringTok{\textquotesingle{}用户 ID\textquotesingle{}}\NormalTok{,}
\NormalTok{  theme\_color }\DataTypeTok{VARCHAR}\NormalTok{(}\DecValTok{20}\NormalTok{) }\KeywordTok{NOT} \KeywordTok{NULL} \KeywordTok{DEFAULT} \StringTok{\textquotesingle{}\#409EFF\textquotesingle{}} \KeywordTok{COMMENT} \StringTok{\textquotesingle{}主题色\textquotesingle{}}\NormalTok{,}
\NormalTok{  background\_mode ENUM(}\StringTok{\textquotesingle{}particle\textquotesingle{}}\NormalTok{, }\StringTok{\textquotesingle{}image\textquotesingle{}}\NormalTok{, }\StringTok{\textquotesingle{}gradient\textquotesingle{}}\NormalTok{) }\KeywordTok{NOT} \KeywordTok{NULL} \KeywordTok{DEFAULT} \StringTok{\textquotesingle{}particle\textquotesingle{}} \KeywordTok{COMMENT} \StringTok{\textquotesingle{}背景类型\textquotesingle{}}\NormalTok{,}
\NormalTok{  background\_asset\_id BIGINT }\KeywordTok{NULL} \KeywordTok{COMMENT} \StringTok{\textquotesingle{}背景资源 ID\textquotesingle{}}\NormalTok{,}
\NormalTok{  layout\_json JSON }\KeywordTok{NULL} \KeywordTok{COMMENT} \StringTok{\textquotesingle{}工作台布局配置，保存主题模式和分页大小等可扩展偏好\textquotesingle{}}\NormalTok{,}
\NormalTok{  confirm\_preferences\_json JSON }\KeywordTok{NULL} \KeywordTok{COMMENT} \StringTok{\textquotesingle{}二次确认偏好\textquotesingle{}}\NormalTok{,}
\NormalTok{  created\_at }\DataTypeTok{TIMESTAMP} \KeywordTok{NOT} \KeywordTok{NULL} \KeywordTok{DEFAULT} \FunctionTok{CURRENT\_TIMESTAMP} \KeywordTok{COMMENT} \StringTok{\textquotesingle{}创建时间\textquotesingle{}}\NormalTok{,}
\NormalTok{  updated\_at }\DataTypeTok{TIMESTAMP} \KeywordTok{NOT} \KeywordTok{NULL} \KeywordTok{DEFAULT} \FunctionTok{CURRENT\_TIMESTAMP} \KeywordTok{ON} \KeywordTok{UPDATE} \FunctionTok{CURRENT\_TIMESTAMP} \KeywordTok{COMMENT} \StringTok{\textquotesingle{}更新时间\textquotesingle{}}\NormalTok{,}

  \KeywordTok{CONSTRAINT}\NormalTok{ fk\_preferences\_user}
    \KeywordTok{FOREIGN} \KeywordTok{KEY}\NormalTok{ (user\_id) }\KeywordTok{REFERENCES}\NormalTok{ users(}\KeywordTok{id}\NormalTok{)}
    \KeywordTok{ON} \KeywordTok{DELETE} \KeywordTok{CASCADE}
\NormalTok{) }\KeywordTok{COMMENT}\OperatorTok{=}\StringTok{\textquotesingle{}用户偏好表\textquotesingle{}}\NormalTok{;}
\end{Highlighting}
\end{Shaded}

说明：

\begin{itemize}
\tightlist
\item
  删除数据库、删除表、删除字段属于强制确认，不由偏好关闭。
\item
  批量插入、批量修改、批量删除、级联删除提示由偏好控制。
\item
  V7 起 \texttt{layout\_json} 增加 \texttt{dashboardLayout}，用于保存
  Dashboard 卡片显示、隐藏、排序和卡片级配置。
\end{itemize}

\hypertarget{user_assets-ux7528ux6237ux8d44ux6e90ux5143ux4fe1ux606fux8868}{%
\subsection{7. user\_assets
用户资源元信息表}\label{user_assets-ux7528ux6237ux8d44ux6e90ux5143ux4fe1ux606fux8868}}

\begin{Shaded}
\begin{Highlighting}[]
\KeywordTok{CREATE} \KeywordTok{TABLE} \ControlFlowTok{IF} \KeywordTok{NOT} \KeywordTok{EXISTS}\NormalTok{ user\_assets (}
  \KeywordTok{id}\NormalTok{ BIGINT AUTO\_INCREMENT }\KeywordTok{PRIMARY} \KeywordTok{KEY} \KeywordTok{COMMENT} \StringTok{\textquotesingle{}资源 ID\textquotesingle{}}\NormalTok{,}
\NormalTok{  user\_id }\DataTypeTok{INT} \KeywordTok{NOT} \KeywordTok{NULL} \KeywordTok{COMMENT} \StringTok{\textquotesingle{}用户 ID\textquotesingle{}}\NormalTok{,}
\NormalTok{  asset\_type ENUM(}\StringTok{\textquotesingle{}avatar\textquotesingle{}}\NormalTok{, }\StringTok{\textquotesingle{}background\textquotesingle{}}\NormalTok{, }\StringTok{\textquotesingle{}other\textquotesingle{}}\NormalTok{) }\KeywordTok{NOT} \KeywordTok{NULL} \KeywordTok{COMMENT} \StringTok{\textquotesingle{}资源类型\textquotesingle{}}\NormalTok{,}
\NormalTok{  storage\_type }\DataTypeTok{VARCHAR}\NormalTok{(}\DecValTok{30}\NormalTok{) }\KeywordTok{NOT} \KeywordTok{NULL} \KeywordTok{DEFAULT} \StringTok{\textquotesingle{}local\textquotesingle{}} \KeywordTok{COMMENT} \StringTok{\textquotesingle{}存储类型\textquotesingle{}}\NormalTok{,}
\NormalTok{  storage\_key }\DataTypeTok{VARCHAR}\NormalTok{(}\DecValTok{500}\NormalTok{) }\KeywordTok{NOT} \KeywordTok{NULL} \KeywordTok{COMMENT} \StringTok{\textquotesingle{}存储键\textquotesingle{}}\NormalTok{,}
\NormalTok{  original\_name }\DataTypeTok{VARCHAR}\NormalTok{(}\DecValTok{255}\NormalTok{) }\KeywordTok{NULL} \KeywordTok{COMMENT} \StringTok{\textquotesingle{}原始文件名\textquotesingle{}}\NormalTok{,}
\NormalTok{  mime\_type }\DataTypeTok{VARCHAR}\NormalTok{(}\DecValTok{100}\NormalTok{) }\KeywordTok{NOT} \KeywordTok{NULL} \KeywordTok{COMMENT} \StringTok{\textquotesingle{}MIME 类型\textquotesingle{}}\NormalTok{,}
\NormalTok{  file\_size BIGINT UNSIGNED }\KeywordTok{NOT} \KeywordTok{NULL} \KeywordTok{COMMENT} \StringTok{\textquotesingle{}文件大小，字节\textquotesingle{}}\NormalTok{,}
\NormalTok{  checksum }\DataTypeTok{VARCHAR}\NormalTok{(}\DecValTok{128}\NormalTok{) }\KeywordTok{NOT} \KeywordTok{NULL} \KeywordTok{COMMENT} \StringTok{\textquotesingle{}文件校验和\textquotesingle{}}\NormalTok{,}
\NormalTok{  status ENUM(}\StringTok{\textquotesingle{}active\textquotesingle{}}\NormalTok{, }\StringTok{\textquotesingle{}deleted\textquotesingle{}}\NormalTok{) }\KeywordTok{NOT} \KeywordTok{NULL} \KeywordTok{DEFAULT} \StringTok{\textquotesingle{}active\textquotesingle{}} \KeywordTok{COMMENT} \StringTok{\textquotesingle{}状态\textquotesingle{}}\NormalTok{,}
\NormalTok{  created\_at }\DataTypeTok{TIMESTAMP} \KeywordTok{NOT} \KeywordTok{NULL} \KeywordTok{DEFAULT} \FunctionTok{CURRENT\_TIMESTAMP} \KeywordTok{COMMENT} \StringTok{\textquotesingle{}创建时间\textquotesingle{}}\NormalTok{,}
\NormalTok{  deleted\_at DATETIME }\KeywordTok{NULL} \KeywordTok{COMMENT} \StringTok{\textquotesingle{}删除时间\textquotesingle{}}\NormalTok{,}

  \KeywordTok{CONSTRAINT}\NormalTok{ fk\_assets\_user}
    \KeywordTok{FOREIGN} \KeywordTok{KEY}\NormalTok{ (user\_id) }\KeywordTok{REFERENCES}\NormalTok{ users(}\KeywordTok{id}\NormalTok{)}
    \KeywordTok{ON} \KeywordTok{DELETE} \KeywordTok{CASCADE}\NormalTok{,}
  \KeywordTok{KEY}\NormalTok{ idx\_assets\_user\_type (user\_id, asset\_type),}
  \KeywordTok{KEY}\NormalTok{ idx\_assets\_user\_time (user\_id, created\_at),}
  \KeywordTok{KEY}\NormalTok{ idx\_assets\_status (status)}
\NormalTok{) }\KeywordTok{COMMENT}\OperatorTok{=}\StringTok{\textquotesingle{}用户资源元信息表\textquotesingle{}}\NormalTok{;}
\end{Highlighting}
\end{Shaded}

说明：

\begin{itemize}
\tightlist
\item
  当前存储类型固定写入 \texttt{local}。
\item
  \texttt{storage\_key} 保存服务端生成的资源键。
\item
  数据库不保存图片二进制。
\end{itemize}

\hypertarget{refresh_tokens-ux4f1aux8bddux5237ux65b0ux8868}{%
\subsection{8. refresh\_tokens
会话刷新表}\label{refresh_tokens-ux4f1aux8bddux5237ux65b0ux8868}}

\begin{Shaded}
\begin{Highlighting}[]
\KeywordTok{CREATE} \KeywordTok{TABLE} \ControlFlowTok{IF} \KeywordTok{NOT} \KeywordTok{EXISTS}\NormalTok{ refresh\_tokens (}
  \KeywordTok{id}\NormalTok{ BIGINT AUTO\_INCREMENT }\KeywordTok{PRIMARY} \KeywordTok{KEY} \KeywordTok{COMMENT} \StringTok{\textquotesingle{}Refresh Token ID\textquotesingle{}}\NormalTok{,}
\NormalTok{  user\_id }\DataTypeTok{INT} \KeywordTok{NOT} \KeywordTok{NULL} \KeywordTok{COMMENT} \StringTok{\textquotesingle{}用户 ID\textquotesingle{}}\NormalTok{,}
\NormalTok{  token\_hash }\DataTypeTok{CHAR}\NormalTok{(}\DecValTok{64}\NormalTok{) }\KeywordTok{NOT} \KeywordTok{NULL} \KeywordTok{COMMENT} \StringTok{\textquotesingle{}Refresh Token SHA{-}256 哈希\textquotesingle{}}\NormalTok{,}
\NormalTok{  token\_family }\DataTypeTok{CHAR}\NormalTok{(}\DecValTok{36}\NormalTok{) }\KeywordTok{NOT} \KeywordTok{NULL} \KeywordTok{COMMENT} \StringTok{\textquotesingle{}Token 家族 ID\textquotesingle{}}\NormalTok{,}
\NormalTok{  expires\_at }\DataTypeTok{TIMESTAMP} \KeywordTok{NOT} \KeywordTok{NULL} \KeywordTok{COMMENT} \StringTok{\textquotesingle{}过期时间\textquotesingle{}}\NormalTok{,}
\NormalTok{  revoked\_at }\DataTypeTok{TIMESTAMP} \KeywordTok{NULL} \KeywordTok{DEFAULT} \KeywordTok{NULL} \KeywordTok{COMMENT} \StringTok{\textquotesingle{}撤销时间\textquotesingle{}}\NormalTok{,}
\NormalTok{  created\_at }\DataTypeTok{TIMESTAMP} \KeywordTok{NOT} \KeywordTok{NULL} \KeywordTok{DEFAULT} \FunctionTok{CURRENT\_TIMESTAMP} \KeywordTok{COMMENT} \StringTok{\textquotesingle{}创建时间\textquotesingle{}}\NormalTok{,}

  \KeywordTok{CONSTRAINT}\NormalTok{ fk\_refresh\_user}
    \KeywordTok{FOREIGN} \KeywordTok{KEY}\NormalTok{ (user\_id) }\KeywordTok{REFERENCES}\NormalTok{ users(}\KeywordTok{id}\NormalTok{)}
    \KeywordTok{ON} \KeywordTok{DELETE} \KeywordTok{CASCADE}\NormalTok{,}
  \KeywordTok{UNIQUE} \KeywordTok{KEY}\NormalTok{ idx\_refresh\_token\_hash (token\_hash),}
  \KeywordTok{KEY}\NormalTok{ idx\_refresh\_user\_family (user\_id, token\_family),}
  \KeywordTok{KEY}\NormalTok{ idx\_refresh\_user\_expires (user\_id, expires\_at)}
\NormalTok{) }\KeywordTok{COMMENT}\OperatorTok{=}\StringTok{\textquotesingle{}Refresh Token 会话表\textquotesingle{}}\NormalTok{;}
\end{Highlighting}
\end{Shaded}

\hypertarget{user_databases-ux7528ux6237ux6570ux636eux5e93ux5f52ux5c5eux8868}{%
\subsection{9. user\_databases
用户数据库归属表}\label{user_databases-ux7528ux6237ux6570ux636eux5e93ux5f52ux5c5eux8868}}

\begin{Shaded}
\begin{Highlighting}[]
\KeywordTok{CREATE} \KeywordTok{TABLE} \ControlFlowTok{IF} \KeywordTok{NOT} \KeywordTok{EXISTS}\NormalTok{ user\_databases (}
  \KeywordTok{id}\NormalTok{ BIGINT AUTO\_INCREMENT }\KeywordTok{PRIMARY} \KeywordTok{KEY} \KeywordTok{COMMENT} \StringTok{\textquotesingle{}用户数据库元数据 ID\textquotesingle{}}\NormalTok{,}
\NormalTok{  user\_id }\DataTypeTok{INT} \KeywordTok{NOT} \KeywordTok{NULL} \KeywordTok{COMMENT} \StringTok{\textquotesingle{}用户 ID\textquotesingle{}}\NormalTok{,}
\NormalTok{  display\_name }\DataTypeTok{VARCHAR}\NormalTok{(}\DecValTok{64}\NormalTok{) }\KeywordTok{NOT} \KeywordTok{NULL} \KeywordTok{COMMENT} \StringTok{\textquotesingle{}用户可见数据库名\textquotesingle{}}\NormalTok{,}
\NormalTok{  schema\_name }\DataTypeTok{VARCHAR}\NormalTok{(}\DecValTok{128}\NormalTok{) }\KeywordTok{NOT} \KeywordTok{NULL} \KeywordTok{COMMENT} \StringTok{\textquotesingle{}实际 MySQL schema 名称\textquotesingle{}}\NormalTok{,}
\NormalTok{  status ENUM(}\StringTok{\textquotesingle{}active\textquotesingle{}}\NormalTok{, }\StringTok{\textquotesingle{}deleting\textquotesingle{}}\NormalTok{, }\StringTok{\textquotesingle{}deleted\textquotesingle{}}\NormalTok{) }\KeywordTok{NOT} \KeywordTok{NULL} \KeywordTok{DEFAULT} \StringTok{\textquotesingle{}active\textquotesingle{}} \KeywordTok{COMMENT} \StringTok{\textquotesingle{}数据库状态\textquotesingle{}}\NormalTok{,}
\NormalTok{  table\_count }\DataTypeTok{INT} \KeywordTok{NOT} \KeywordTok{NULL} \KeywordTok{DEFAULT} \DecValTok{0} \KeywordTok{COMMENT} \StringTok{\textquotesingle{}表数量快照\textquotesingle{}}\NormalTok{,}
\NormalTok{  size\_bytes BIGINT }\KeywordTok{NOT} \KeywordTok{NULL} \KeywordTok{DEFAULT} \DecValTok{0} \KeywordTok{COMMENT} \StringTok{\textquotesingle{}容量快照\textquotesingle{}}\NormalTok{,}
\NormalTok{  created\_at }\DataTypeTok{TIMESTAMP} \KeywordTok{NOT} \KeywordTok{NULL} \KeywordTok{DEFAULT} \FunctionTok{CURRENT\_TIMESTAMP} \KeywordTok{COMMENT} \StringTok{\textquotesingle{}创建时间\textquotesingle{}}\NormalTok{,}
\NormalTok{  updated\_at }\DataTypeTok{TIMESTAMP} \KeywordTok{NOT} \KeywordTok{NULL} \KeywordTok{DEFAULT} \FunctionTok{CURRENT\_TIMESTAMP} \KeywordTok{ON} \KeywordTok{UPDATE} \FunctionTok{CURRENT\_TIMESTAMP} \KeywordTok{COMMENT} \StringTok{\textquotesingle{}更新时间\textquotesingle{}}\NormalTok{,}
\NormalTok{  deleted\_at }\DataTypeTok{TIMESTAMP} \KeywordTok{NULL} \KeywordTok{DEFAULT} \KeywordTok{NULL} \KeywordTok{COMMENT} \StringTok{\textquotesingle{}删除时间\textquotesingle{}}\NormalTok{,}

  \KeywordTok{CONSTRAINT}\NormalTok{ fk\_user\_databases\_user}
    \KeywordTok{FOREIGN} \KeywordTok{KEY}\NormalTok{ (user\_id) }\KeywordTok{REFERENCES}\NormalTok{ users(}\KeywordTok{id}\NormalTok{)}
    \KeywordTok{ON} \KeywordTok{DELETE} \KeywordTok{CASCADE}\NormalTok{,}
  \KeywordTok{UNIQUE} \KeywordTok{KEY}\NormalTok{ uk\_user\_databases\_schema\_name (schema\_name),}
  \KeywordTok{UNIQUE} \KeywordTok{KEY}\NormalTok{ uk\_user\_databases\_user\_display\_status (user\_id, display\_name, status),}
  \KeywordTok{KEY}\NormalTok{ idx\_user\_databases\_user\_status\_time (user\_id, status, created\_at)}
\NormalTok{) }\KeywordTok{COMMENT}\OperatorTok{=}\StringTok{\textquotesingle{}用户数据库元数据表\textquotesingle{}}\NormalTok{;}
\end{Highlighting}
\end{Shaded}

物理库命名格式：

\begin{Shaded}
\begin{Highlighting}[]
\NormalTok{u\{userId\}\_\{safeDisplayName\}\_\{randomSuffix\}}
\end{Highlighting}
\end{Shaded}

说明：

\begin{itemize}
\tightlist
\item
  \texttt{schema\_name}
  只允许后端内部使用，前端和接口响应不返回真实物理库名。
\item
  \texttt{status=\textquotesingle{}deleting\textquotesingle{}}
  表示正在执行物理库删除，用于防止删除失败时状态与 MySQL schema 不一致。
\item
  \texttt{table\_count} 和 \texttt{size\_bytes}
  是快照字段；列表接口优先从 \texttt{information\_schema.tables}
  读取实时值。
\item
  服务启动时会检查并创建
  \texttt{user\_databases}，避免已有开发库漏执行新版 \texttt{init.sql}。
\end{itemize}

\hypertarget{transaction_sessions-ux4e8bux52a1ux4f1aux8bddux8868}{%
\subsection{10. transaction\_sessions
事务会话表}\label{transaction_sessions-ux4e8bux52a1ux4f1aux8bddux8868}}

\begin{Shaded}
\begin{Highlighting}[]
\KeywordTok{CREATE} \KeywordTok{TABLE} \ControlFlowTok{IF} \KeywordTok{NOT} \KeywordTok{EXISTS}\NormalTok{ transaction\_sessions (}
  \KeywordTok{id}\NormalTok{ BIGINT UNSIGNED AUTO\_INCREMENT }\KeywordTok{PRIMARY} \KeywordTok{KEY} \KeywordTok{COMMENT} \StringTok{\textquotesingle{}事务会话 ID\textquotesingle{}}\NormalTok{,}
\NormalTok{  transaction\_key }\DataTypeTok{VARCHAR}\NormalTok{(}\DecValTok{64}\NormalTok{) }\KeywordTok{NOT} \KeywordTok{NULL} \KeywordTok{UNIQUE} \KeywordTok{COMMENT} \StringTok{\textquotesingle{}对外事务 ID\textquotesingle{}}\NormalTok{,}
\NormalTok{  user\_id }\DataTypeTok{INT} \KeywordTok{NOT} \KeywordTok{NULL} \KeywordTok{COMMENT} \StringTok{\textquotesingle{}用户 ID\textquotesingle{}}\NormalTok{,}
\NormalTok{  database\_id BIGINT }\KeywordTok{NOT} \KeywordTok{NULL} \KeywordTok{COMMENT} \StringTok{\textquotesingle{}平台数据库 ID\textquotesingle{}}\NormalTok{,}
  \KeywordTok{isolation\_level}\NormalTok{ ENUM(}\StringTok{\textquotesingle{}READ COMMITTED\textquotesingle{}}\NormalTok{, }\StringTok{\textquotesingle{}REPEATABLE READ\textquotesingle{}}\NormalTok{, }\StringTok{\textquotesingle{}SERIALIZABLE\textquotesingle{}}\NormalTok{) }\KeywordTok{NOT} \KeywordTok{NULL} \KeywordTok{DEFAULT} \StringTok{\textquotesingle{}REPEATABLE READ\textquotesingle{}} \KeywordTok{COMMENT} \StringTok{\textquotesingle{}隔离级别\textquotesingle{}}\NormalTok{,}
\NormalTok{  status ENUM(}\StringTok{\textquotesingle{}active\textquotesingle{}}\NormalTok{, }\StringTok{\textquotesingle{}committed\textquotesingle{}}\NormalTok{, }\StringTok{\textquotesingle{}rolled\_back\textquotesingle{}}\NormalTok{, }\StringTok{\textquotesingle{}expired\textquotesingle{}}\NormalTok{) }\KeywordTok{NOT} \KeywordTok{NULL} \KeywordTok{DEFAULT} \StringTok{\textquotesingle{}active\textquotesingle{}} \KeywordTok{COMMENT} \StringTok{\textquotesingle{}事务状态\textquotesingle{}}\NormalTok{,}
\NormalTok{  started\_at }\DataTypeTok{TIMESTAMP} \KeywordTok{NOT} \KeywordTok{NULL} \KeywordTok{DEFAULT} \FunctionTok{CURRENT\_TIMESTAMP} \KeywordTok{COMMENT} \StringTok{\textquotesingle{}开始时间\textquotesingle{}}\NormalTok{,}
\NormalTok{  expires\_at DATETIME }\KeywordTok{NOT} \KeywordTok{NULL} \KeywordTok{COMMENT} \StringTok{\textquotesingle{}过期时间\textquotesingle{}}\NormalTok{,}
\NormalTok{  finished\_at DATETIME }\KeywordTok{NULL} \KeywordTok{COMMENT} \StringTok{\textquotesingle{}结束时间\textquotesingle{}}\NormalTok{,}

  \KeywordTok{INDEX}\NormalTok{ idx\_transaction\_user\_status (user\_id, status),}
  \KeywordTok{INDEX}\NormalTok{ idx\_transaction\_database\_status (database\_id, status),}
  \KeywordTok{INDEX}\NormalTok{ idx\_transaction\_expires (expires\_at),}

  \KeywordTok{CONSTRAINT}\NormalTok{ fk\_transaction\_user}
    \KeywordTok{FOREIGN} \KeywordTok{KEY}\NormalTok{ (user\_id) }\KeywordTok{REFERENCES}\NormalTok{ users(}\KeywordTok{id}\NormalTok{),}
  \KeywordTok{CONSTRAINT}\NormalTok{ fk\_transaction\_database}
    \KeywordTok{FOREIGN} \KeywordTok{KEY}\NormalTok{ (database\_id) }\KeywordTok{REFERENCES}\NormalTok{ user\_databases(}\KeywordTok{id}\NormalTok{)}
\NormalTok{) }\KeywordTok{COMMENT}\OperatorTok{=}\StringTok{\textquotesingle{}事务会话表\textquotesingle{}}\NormalTok{;}
\end{Highlighting}
\end{Shaded}

说明：

\begin{itemize}
\tightlist
\item
  \texttt{transaction\_key} 是前端使用的事务 ID。
\item
  实际数据库连接由后端内存连接池管理。
\item
  表中记录用于追踪事务生命周期。
\end{itemize}

\hypertarget{audit_logs-ux64cdux4f5cux65e5ux5fd7ux8868}{%
\subsection{11. audit\_logs
操作日志表}\label{audit_logs-ux64cdux4f5cux65e5ux5fd7ux8868}}

\begin{Shaded}
\begin{Highlighting}[]
\KeywordTok{CREATE} \KeywordTok{TABLE} \ControlFlowTok{IF} \KeywordTok{NOT} \KeywordTok{EXISTS}\NormalTok{ audit\_logs (}
  \KeywordTok{id}\NormalTok{ BIGINT UNSIGNED AUTO\_INCREMENT }\KeywordTok{PRIMARY} \KeywordTok{KEY} \KeywordTok{COMMENT} \StringTok{\textquotesingle{}日志 ID\textquotesingle{}}\NormalTok{,}
\NormalTok{  trace\_id }\DataTypeTok{VARCHAR}\NormalTok{(}\DecValTok{64}\NormalTok{) }\KeywordTok{NOT} \KeywordTok{NULL} \KeywordTok{COMMENT} \StringTok{\textquotesingle{}请求追踪 ID\textquotesingle{}}\NormalTok{,}
\NormalTok{  user\_id }\DataTypeTok{INT} \KeywordTok{NULL} \KeywordTok{COMMENT} \StringTok{\textquotesingle{}操作用户 ID\textquotesingle{}}\NormalTok{,}
\NormalTok{  database\_id BIGINT }\KeywordTok{NULL} \KeywordTok{COMMENT} \StringTok{\textquotesingle{}平台数据库 ID\textquotesingle{}}\NormalTok{,}
\NormalTok{  transaction\_id BIGINT UNSIGNED }\KeywordTok{NULL} \KeywordTok{COMMENT} \StringTok{\textquotesingle{}事务会话 ID\textquotesingle{}}\NormalTok{,}
\NormalTok{  object\_type }\DataTypeTok{VARCHAR}\NormalTok{(}\DecValTok{30}\NormalTok{) }\KeywordTok{NOT} \KeywordTok{NULL} \KeywordTok{COMMENT} \StringTok{\textquotesingle{}对象类型\textquotesingle{}}\NormalTok{,}
\NormalTok{  object\_name }\DataTypeTok{VARCHAR}\NormalTok{(}\DecValTok{128}\NormalTok{) }\KeywordTok{NULL} \KeywordTok{COMMENT} \StringTok{\textquotesingle{}对象名称\textquotesingle{}}\NormalTok{,}
\NormalTok{  action\_type }\DataTypeTok{VARCHAR}\NormalTok{(}\DecValTok{50}\NormalTok{) }\KeywordTok{NOT} \KeywordTok{NULL} \KeywordTok{COMMENT} \StringTok{\textquotesingle{}操作类型\textquotesingle{}}\NormalTok{,}
  \KeywordTok{summary} \DataTypeTok{VARCHAR}\NormalTok{(}\DecValTok{255}\NormalTok{) }\KeywordTok{NOT} \KeywordTok{NULL} \KeywordTok{COMMENT} \StringTok{\textquotesingle{}操作摘要\textquotesingle{}}\NormalTok{,}
\NormalTok{  detail\_json JSON }\KeywordTok{NULL} \KeywordTok{COMMENT} \StringTok{\textquotesingle{}操作详情\textquotesingle{}}\NormalTok{,}
\NormalTok{  success }\DataTypeTok{BOOLEAN} \KeywordTok{NOT} \KeywordTok{NULL} \KeywordTok{DEFAULT} \KeywordTok{TRUE} \KeywordTok{COMMENT} \StringTok{\textquotesingle{}成功状态\textquotesingle{}}\NormalTok{,}
\NormalTok{  error\_code }\DataTypeTok{VARCHAR}\NormalTok{(}\DecValTok{80}\NormalTok{) }\KeywordTok{NULL} \KeywordTok{COMMENT} \StringTok{\textquotesingle{}错误码\textquotesingle{}}\NormalTok{,}
\NormalTok{  error\_message }\DataTypeTok{VARCHAR}\NormalTok{(}\DecValTok{500}\NormalTok{) }\KeywordTok{NULL} \KeywordTok{COMMENT} \StringTok{\textquotesingle{}错误消息\textquotesingle{}}\NormalTok{,}
\NormalTok{  duration\_ms }\DataTypeTok{INT}\NormalTok{ UNSIGNED }\KeywordTok{NULL} \KeywordTok{COMMENT} \StringTok{\textquotesingle{}耗时毫秒\textquotesingle{}}\NormalTok{,}
\NormalTok{  ip\_address }\DataTypeTok{VARCHAR}\NormalTok{(}\DecValTok{45}\NormalTok{) }\KeywordTok{NULL} \KeywordTok{COMMENT} \StringTok{\textquotesingle{}客户端 IP\textquotesingle{}}\NormalTok{,}
\NormalTok{  user\_agent }\DataTypeTok{VARCHAR}\NormalTok{(}\DecValTok{255}\NormalTok{) }\KeywordTok{NULL} \KeywordTok{COMMENT} \StringTok{\textquotesingle{}User{-}Agent\textquotesingle{}}\NormalTok{,}
\NormalTok{  created\_at }\DataTypeTok{TIMESTAMP} \KeywordTok{NOT} \KeywordTok{NULL} \KeywordTok{DEFAULT} \FunctionTok{CURRENT\_TIMESTAMP} \KeywordTok{COMMENT} \StringTok{\textquotesingle{}创建时间\textquotesingle{}}\NormalTok{,}

  \KeywordTok{INDEX}\NormalTok{ idx\_audit\_trace (trace\_id),}
  \KeywordTok{INDEX}\NormalTok{ idx\_audit\_user\_time (user\_id, created\_at),}
  \KeywordTok{INDEX}\NormalTok{ idx\_audit\_database\_time (database\_id, created\_at),}
  \KeywordTok{INDEX}\NormalTok{ idx\_audit\_action\_type (action\_type),}
  \KeywordTok{INDEX}\NormalTok{ idx\_audit\_object\_type (object\_type),}

  \KeywordTok{CONSTRAINT}\NormalTok{ fk\_audit\_user}
    \KeywordTok{FOREIGN} \KeywordTok{KEY}\NormalTok{ (user\_id) }\KeywordTok{REFERENCES}\NormalTok{ users(}\KeywordTok{id}\NormalTok{),}
  \KeywordTok{CONSTRAINT}\NormalTok{ fk\_audit\_database}
    \KeywordTok{FOREIGN} \KeywordTok{KEY}\NormalTok{ (database\_id) }\KeywordTok{REFERENCES}\NormalTok{ user\_databases(}\KeywordTok{id}\NormalTok{),}
  \KeywordTok{CONSTRAINT}\NormalTok{ fk\_audit\_transaction}
    \KeywordTok{FOREIGN} \KeywordTok{KEY}\NormalTok{ (transaction\_id) }\KeywordTok{REFERENCES}\NormalTok{ transaction\_sessions(}\KeywordTok{id}\NormalTok{)}
\NormalTok{) }\KeywordTok{COMMENT}\OperatorTok{=}\StringTok{\textquotesingle{}操作审计日志表\textquotesingle{}}\NormalTok{;}
\end{Highlighting}
\end{Shaded}

\texttt{object\_type} 值：

\begin{Shaded}
\begin{Highlighting}[]
\NormalTok{AUTH}
\NormalTok{USER}
\NormalTok{PREFERENCE}
\NormalTok{DATABASE}
\NormalTok{TABLE}
\NormalTok{COLUMN}
\NormalTok{CONSTRAINT}
\NormalTok{INDEX}
\NormalTok{VIEW}
\NormalTok{ROW}
\NormalTok{QUERY}
\NormalTok{TRANSACTION}
\NormalTok{STATS}
\NormalTok{ASSET}
\end{Highlighting}
\end{Shaded}

\texttt{action\_type} 值：

\begin{Shaded}
\begin{Highlighting}[]
\NormalTok{AUTH\_REGISTER}
\NormalTok{AUTH\_LOGIN}
\NormalTok{AUTH\_REFRESH}
\NormalTok{AUTH\_LOGOUT}
\NormalTok{PREFERENCE\_UPDATE}
\NormalTok{DB\_CREATE}
\NormalTok{DB\_UPDATE}
\NormalTok{DB\_DELETE}
\NormalTok{TABLE\_CREATE}
\NormalTok{TABLE\_ALTER}
\NormalTok{TABLE\_DELETE}
\NormalTok{COLUMN\_CREATE}
\NormalTok{COLUMN\_UPDATE}
\NormalTok{COLUMN\_DELETE}
\NormalTok{CONSTRAINT\_CREATE}
\NormalTok{CONSTRAINT\_DELETE}
\NormalTok{INDEX\_CREATE}
\NormalTok{INDEX\_DELETE}
\NormalTok{VIEW\_CREATE}
\NormalTok{VIEW\_UPDATE}
\NormalTok{VIEW\_DELETE}
\NormalTok{ROW\_SELECT}
\NormalTok{ROW\_INSERT}
\NormalTok{ROW\_UPDATE}
\NormalTok{ROW\_DELETE}
\NormalTok{QUERY\_SELECT}
\NormalTok{TX\_BEGIN}
\NormalTok{TX\_COMMIT}
\NormalTok{TX\_ROLLBACK}
\NormalTok{STATS\_VIEW}
\NormalTok{ASSET\_VIEW}
\end{Highlighting}
\end{Shaded}

\hypertarget{stat_snapshots-ux7edfux8ba1ux5febux7167ux8868}{%
\subsection{12. stat\_snapshots
统计快照表}\label{stat_snapshots-ux7edfux8ba1ux5febux7167ux8868}}

\begin{Shaded}
\begin{Highlighting}[]
\KeywordTok{CREATE} \KeywordTok{TABLE} \ControlFlowTok{IF} \KeywordTok{NOT} \KeywordTok{EXISTS}\NormalTok{ stat\_snapshots (}
  \KeywordTok{id}\NormalTok{ BIGINT UNSIGNED AUTO\_INCREMENT }\KeywordTok{PRIMARY} \KeywordTok{KEY} \KeywordTok{COMMENT} \StringTok{\textquotesingle{}快照 ID\textquotesingle{}}\NormalTok{,}
\NormalTok{  user\_id }\DataTypeTok{INT} \KeywordTok{NOT} \KeywordTok{NULL} \KeywordTok{COMMENT} \StringTok{\textquotesingle{}用户 ID\textquotesingle{}}\NormalTok{,}
\NormalTok{  database\_id BIGINT }\KeywordTok{NULL} \KeywordTok{COMMENT} \StringTok{\textquotesingle{}平台数据库 ID\textquotesingle{}}\NormalTok{,}
\NormalTok{  table\_name }\DataTypeTok{VARCHAR}\NormalTok{(}\DecValTok{64}\NormalTok{) }\KeywordTok{NULL} \KeywordTok{COMMENT} \StringTok{\textquotesingle{}表名\textquotesingle{}}\NormalTok{,}
  \KeywordTok{scope}\NormalTok{ ENUM(}\StringTok{\textquotesingle{}user\textquotesingle{}}\NormalTok{, }\StringTok{\textquotesingle{}database\textquotesingle{}}\NormalTok{, }\StringTok{\textquotesingle{}table\textquotesingle{}}\NormalTok{) }\KeywordTok{NOT} \KeywordTok{NULL} \KeywordTok{COMMENT} \StringTok{\textquotesingle{}统计范围\textquotesingle{}}\NormalTok{,}
\NormalTok{  metric\_json JSON }\KeywordTok{NOT} \KeywordTok{NULL} \KeywordTok{COMMENT} \StringTok{\textquotesingle{}统计指标\textquotesingle{}}\NormalTok{,}
\NormalTok{  captured\_at }\DataTypeTok{TIMESTAMP} \KeywordTok{NOT} \KeywordTok{NULL} \KeywordTok{DEFAULT} \FunctionTok{CURRENT\_TIMESTAMP} \KeywordTok{COMMENT} \StringTok{\textquotesingle{}采集时间\textquotesingle{}}\NormalTok{,}

  \KeywordTok{INDEX}\NormalTok{ idx\_stat\_user\_scope\_time (user\_id, }\KeywordTok{scope}\NormalTok{, captured\_at),}
  \KeywordTok{INDEX}\NormalTok{ idx\_stat\_database\_time (database\_id, captured\_at),}

  \KeywordTok{CONSTRAINT}\NormalTok{ fk\_stat\_user}
    \KeywordTok{FOREIGN} \KeywordTok{KEY}\NormalTok{ (user\_id) }\KeywordTok{REFERENCES}\NormalTok{ users(}\KeywordTok{id}\NormalTok{),}
  \KeywordTok{CONSTRAINT}\NormalTok{ fk\_stat\_database}
    \KeywordTok{FOREIGN} \KeywordTok{KEY}\NormalTok{ (database\_id) }\KeywordTok{REFERENCES}\NormalTok{ user\_databases(}\KeywordTok{id}\NormalTok{)}
\NormalTok{) }\KeywordTok{COMMENT}\OperatorTok{=}\StringTok{\textquotesingle{}统计快照表\textquotesingle{}}\NormalTok{;}
\end{Highlighting}
\end{Shaded}

\hypertarget{mysql-ux5143ux6570ux636eux6765ux6e90}{%
\subsection{13. MySQL
元数据来源}\label{mysql-ux5143ux6570ux636eux6765ux6e90}}

表、字段、索引、约束、视图信息以 MySQL 元数据为准。

表与视图：

\begin{Shaded}
\begin{Highlighting}[]
\KeywordTok{SELECT}\NormalTok{ table\_name, table\_type, table\_rows, data\_length, index\_length, create\_time, update\_time}
\KeywordTok{FROM}\NormalTok{ information\_schema.}\KeywordTok{tables}
\KeywordTok{WHERE}\NormalTok{ table\_schema }\OperatorTok{=}\NormalTok{ ?;}
\end{Highlighting}
\end{Shaded}

字段：

\begin{Shaded}
\begin{Highlighting}[]
\KeywordTok{SELECT}
\NormalTok{  column\_name,}
\NormalTok{  data\_type,}
\NormalTok{  column\_type,}
\NormalTok{  is\_nullable,}
\NormalTok{  column\_key,}
\NormalTok{  column\_default,}
\NormalTok{  extra,}
\NormalTok{  ordinal\_position}
\KeywordTok{FROM}\NormalTok{ information\_schema.}\KeywordTok{columns}
\KeywordTok{WHERE}\NormalTok{ table\_schema }\OperatorTok{=}\NormalTok{ ?}
  \KeywordTok{AND}\NormalTok{ table\_name }\OperatorTok{=}\NormalTok{ ?}
\KeywordTok{ORDER} \KeywordTok{BY}\NormalTok{ ordinal\_position;}
\end{Highlighting}
\end{Shaded}

索引：

\begin{Shaded}
\begin{Highlighting}[]
\KeywordTok{SELECT}
\NormalTok{  index\_name,}
\NormalTok{  non\_unique,}
\NormalTok{  seq\_in\_index,}
\NormalTok{  column\_name,}
\NormalTok{  collation,}
\NormalTok{  index\_type}
\KeywordTok{FROM}\NormalTok{ information\_schema.}\KeywordTok{statistics}
\KeywordTok{WHERE}\NormalTok{ table\_schema }\OperatorTok{=}\NormalTok{ ?}
  \KeywordTok{AND}\NormalTok{ table\_name }\OperatorTok{=}\NormalTok{ ?}
\KeywordTok{ORDER} \KeywordTok{BY}\NormalTok{ index\_name, seq\_in\_index;}
\end{Highlighting}
\end{Shaded}

约束：

\begin{Shaded}
\begin{Highlighting}[]
\KeywordTok{SELECT}
\NormalTok{  constraint\_name,}
\NormalTok{  constraint\_type}
\KeywordTok{FROM}\NormalTok{ information\_schema.table\_constraints}
\KeywordTok{WHERE}\NormalTok{ table\_schema }\OperatorTok{=}\NormalTok{ ?}
  \KeywordTok{AND}\NormalTok{ table\_name }\OperatorTok{=}\NormalTok{ ?;}
\end{Highlighting}
\end{Shaded}

外键：

\begin{Shaded}
\begin{Highlighting}[]
\KeywordTok{SELECT}
\NormalTok{  constraint\_name,}
\NormalTok{  column\_name,}
\NormalTok{  referenced\_table\_name,}
\NormalTok{  referenced\_column\_name}
\KeywordTok{FROM}\NormalTok{ information\_schema.key\_column\_usage}
\KeywordTok{WHERE}\NormalTok{ table\_schema }\OperatorTok{=}\NormalTok{ ?}
  \KeywordTok{AND}\NormalTok{ table\_name }\OperatorTok{=}\NormalTok{ ?}
  \KeywordTok{AND}\NormalTok{ referenced\_table\_name }\KeywordTok{IS} \KeywordTok{NOT} \KeywordTok{NULL}\NormalTok{;}
\end{Highlighting}
\end{Shaded}

视图：

\begin{Shaded}
\begin{Highlighting}[]
\KeywordTok{SELECT}\NormalTok{ table\_name, view\_definition, check\_option, is\_updatable}
\KeywordTok{FROM}\NormalTok{ information\_schema.views}
\KeywordTok{WHERE}\NormalTok{ table\_schema }\OperatorTok{=}\NormalTok{ ?}
  \KeywordTok{AND}\NormalTok{ table\_name }\OperatorTok{=}\NormalTok{ ?;}
\end{Highlighting}
\end{Shaded}

\hypertarget{ux652fux6301ux7684ux6570ux636eux7c7bux578b}{%
\subsection{14.
支持的数据类型}\label{ux652fux6301ux7684ux6570ux636eux7c7bux578b}}

GUI 字段类型清单：

\begin{Shaded}
\begin{Highlighting}[]
\NormalTok{TINYINT}
\NormalTok{SMALLINT}
\NormalTok{INT}
\NormalTok{BIGINT}
\NormalTok{DECIMAL}
\NormalTok{FLOAT}
\NormalTok{DOUBLE}
\NormalTok{CHAR}
\NormalTok{VARCHAR}
\NormalTok{TEXT}
\NormalTok{MEDIUMTEXT}
\NormalTok{LONGTEXT}
\NormalTok{DATE}
\NormalTok{TIME}
\NormalTok{DATETIME}
\NormalTok{TIMESTAMP}
\NormalTok{BOOLEAN}
\NormalTok{ENUM}
\NormalTok{JSON}
\NormalTok{BLOB}
\NormalTok{MEDIUMBLOB}
\NormalTok{LONGBLOB}
\end{Highlighting}
\end{Shaded}

字段属性：

\begin{itemize}
\tightlist
\item
  unsigned
\item
  length
\item
  precision
\item
  scale
\item
  nullable
\item
  defaultValue
\item
  autoIncrement
\item
  comment
\item
  enumValues：枚举式有限取值，落地时优先转换为 MySQL ENUM 或受控 CHECK
  约束。
\item
  checkExpression：检查约束表达式，必须由 GUI
  结构化配置并由后端白名单生成 SQL。
\end{itemize}

\hypertarget{ux7d22ux5f15ux8bbeux8ba1}{%
\subsection{15. 索引设计}\label{ux7d22ux5f15ux8bbeux8ba1}}

\hypertarget{data-idx-001-ux7d22ux5f15ux8bbeux8ba1ux539fux5219}{%
\subsubsection{DATA-IDX-001
索引设计原则}\label{data-idx-001-ux7d22ux5f15ux8bbeux8ba1ux539fux5219}}

系统库索引必须服务于高频查询、归属校验、唯一性约束和审计检索。

索引设计原则：

\begin{itemize}
\tightlist
\item
  主键使用自增 ID 或稳定唯一 ID。
\item
  外键字段必须建立索引。
\item
  高频登录字段必须唯一索引。
\item
  高频按用户查询的表必须包含 \texttt{user\_id} 前缀索引。
\item
  高频时间列表查询使用组合索引，例如 \texttt{(user\_id,\ created\_at)}。
\item
  不为低选择性布尔字段单独建索引，除非与高选择性字段组合。
\item
  JSON 字段默认不建索引；若后续需要查询 JSON
  内部字段，再考虑生成列索引。
\end{itemize}

\hypertarget{data-idx-002-ux7cfbux7edfux5e93ux7d22ux5f15ux6e05ux5355}{%
\subsubsection{DATA-IDX-002
系统库索引清单}\label{data-idx-002-ux7cfbux7edfux5e93ux7d22ux5f15ux6e05ux5355}}

为避免长索引名在 PDF
表格中挤压列宽，本清单按表分组维护；每条记录仍保持``索引、类型、目的''三项信息。

\begin{itemize}
\tightlist
\item
  \texttt{users}

  \begin{itemize}
  \tightlist
  \item
    \texttt{uk\_users\_account\_no(account\_no)}：唯一索引；登录账号唯一和登录查询。
  \item
    \texttt{idx\_users\_status(status)}：普通索引；后台筛选用户状态，当前阶段可选。
  \end{itemize}
\item
  \texttt{user\_profiles}

  \begin{itemize}
  \tightlist
  \item
    \texttt{pk\_user\_profiles(user\_id)}：主键/唯一；一个用户一份资料。
  \end{itemize}
\item
  \texttt{user\_preferences}

  \begin{itemize}
  \tightlist
  \item
    \texttt{pk\_user\_preferences(user\_id)}：主键/唯一；一个用户一份偏好。
  \end{itemize}
\item
  \texttt{user\_assets}

  \begin{itemize}
  \tightlist
  \item
    \texttt{idx\_assets\_user\_type(user\_id,\ asset\_type)}：普通索引；查询用户头像、背景等资源。
  \item
    \texttt{idx\_assets\_user\_time(user\_id,\ created\_at)}：普通索引；用户资源列表。
  \end{itemize}
\item
  \texttt{refresh\_tokens}

  \begin{itemize}
  \tightlist
  \item
    \texttt{idx\_refresh\_token\_hash(token\_hash)}：唯一索引；快速校验
    Refresh Token。
  \item
    \texttt{idx\_refresh\_user\_family(user\_id,\ token\_family)}：普通索引；Token
    家族撤销。
  \item
    \texttt{idx\_refresh\_user\_expires(user\_id,\ expires\_at)}：普通索引；清理过期
    Token。
  \end{itemize}
\item
  \texttt{user\_databases}

  \begin{itemize}
  \tightlist
  \item
    \texttt{uk\_user\_databases\_schema\_name(schema\_name)}：唯一索引；真实物理库名唯一。
  \item
    \texttt{uk\_user\_databases\_user\_display\_status(user\_id,\ display\_name,\ status)}：唯一索引；同一用户同一状态下数据库显示名唯一。
  \item
    \texttt{idx\_user\_databases\_user\_status\_time(user\_id,\ status,\ created\_at)}：普通索引；查询当前用户有效数据库并按创建时间排序。
  \end{itemize}
\item
  \texttt{transaction\_sessions}

  \begin{itemize}
  \tightlist
  \item
    \texttt{idx\_tx\_user\_status(user\_id,\ status)}：普通索引；查询用户活跃事务。
  \item
    \texttt{idx\_tx\_expires\_at(expires\_at)}：普通索引；超时事务扫描。
  \end{itemize}
\item
  \texttt{audit\_logs}

  \begin{itemize}
  \tightlist
  \item
    \texttt{idx\_audit\_user\_time(user\_id,\ created\_at)}：普通索引；用户操作日志分页。
  \item
    \texttt{idx\_audit\_database\_time(database\_id,\ created\_at)}：普通索引；数据库级操作日志。
  \item
    \texttt{idx\_audit\_action\_time(action\_type,\ created\_at)}：普通索引；操作类型统计。
  \item
    \texttt{idx\_audit\_trace\_id(trace\_id)}：普通索引；根据 traceId
    排查问题。
  \end{itemize}
\item
  \texttt{stat\_snapshots}

  \begin{itemize}
  \tightlist
  \item
    \texttt{idx\_stat\_user\_scope\_time(user\_id,\ scope,\ captured\_at)}：普通索引；用户统计趋势。
  \item
    \texttt{idx\_stat\_database\_time(database\_id,\ captured\_at)}：普通索引；数据库统计趋势。
  \end{itemize}
\end{itemize}

\hypertarget{data-idx-003-ux7528ux6237ux7269ux7406ux5e93ux7d22ux5f15ux7b56ux7565}{%
\subsubsection{DATA-IDX-003
用户物理库索引策略}\label{data-idx-003-ux7528ux6237ux7269ux7406ux5e93ux7d22ux5f15ux7b56ux7565}}

用户创建表时，索引由 GUI 显式配置。系统不自动为所有字段创建索引。

默认策略：

\begin{itemize}
\tightlist
\item
  主键必须创建主键索引。
\item
  唯一约束必须创建唯一索引。
\item
  外键字段建议创建普通索引；若 MySQL 自动创建则复用。
\item
  用户显式创建普通索引或唯一索引时，字段必须存在。
\item
  单表索引数量默认最多 \texttt{16} 个。
\item
  复合索引字段顺序由用户在 GUI 中配置。
\end{itemize}

\hypertarget{data-idx-004-ux7d22ux5f15ux53d8ux66f4ux5b89ux5168}{%
\subsubsection{DATA-IDX-004
索引变更安全}\label{data-idx-004-ux7d22ux5f15ux53d8ux66f4ux5b89ux5168}}

创建或删除索引前必须：

\begin{itemize}
\tightlist
\item
  校验数据库归属。
\item
  校验表存在。
\item
  校验字段存在。
\item
  校验索引名符合标识符规则。
\item
  校验单表索引数量上限。
\item
  对删除唯一索引、删除外键相关索引等高风险操作要求二次确认。
\item
  写入审计日志。
\end{itemize}

\hypertarget{data-idx-005-ux7d22ux5f15ux7edfux8ba1ux6765ux6e90}{%
\subsubsection{DATA-IDX-005
索引统计来源}\label{data-idx-005-ux7d22ux5f15ux7edfux8ba1ux6765ux6e90}}

索引元数据以 MySQL 为准：

\begin{Shaded}
\begin{Highlighting}[]
\KeywordTok{SELECT}
\NormalTok{  index\_name,}
\NormalTok{  non\_unique,}
\NormalTok{  column\_name,}
\NormalTok{  seq\_in\_index,}
\NormalTok{  collation,}
\NormalTok{  index\_type}
\KeywordTok{FROM}\NormalTok{ information\_schema.}\KeywordTok{statistics}
\KeywordTok{WHERE}\NormalTok{ table\_schema }\OperatorTok{=}\NormalTok{ ?}
  \KeywordTok{AND}\NormalTok{ table\_name }\OperatorTok{=}\NormalTok{ ?}
\KeywordTok{ORDER} \KeywordTok{BY}\NormalTok{ index\_name, seq\_in\_index;}
\end{Highlighting}
\end{Shaded}

\hypertarget{ux5bb9ux91cfux9650ux5236}{%
\subsection{16. 容量限制}\label{ux5bb9ux91cfux9650ux5236}}

\hypertarget{data-limit-001-ux9650ux5236ux6765ux6e90}{%
\subsubsection{DATA-LIMIT-001
限制来源}\label{data-limit-001-ux9650ux5236ux6765ux6e90}}

容量限制来自服务端配置。系统库默认不保存限制值，避免把部署策略和业务数据耦合。

如果后续需要不同用户套餐，可以新增独立表：

\begin{Shaded}
\begin{Highlighting}[]
\NormalTok{plans}
\NormalTok{user\_plan\_assignments}
\end{Highlighting}
\end{Shaded}

当前阶段不设计套餐表，统一使用服务端默认限制。

\hypertarget{data-limit-002-ux9ed8ux8ba4ux5bb9ux91cfux9650ux5236}{%
\subsubsection{DATA-LIMIT-002
默认容量限制}\label{data-limit-002-ux9ed8ux8ba4ux5bb9ux91cfux9650ux5236}}

\begin{longtable}[]{@{}lll@{}}
\toprule
限制项 & 默认值 & 校验依据 \\
\midrule
\endhead
单用户数据库数量 & 10 & \texttt{user\_databases} \\
单数据库表数量 & 100 & \texttt{information\_schema.tables} \\
单数据库视图数量 & 50 & \texttt{information\_schema.views} \\
单表字段数量 & 80 & \texttt{information\_schema.columns} \\
单表索引数量 & 16 & \texttt{information\_schema.statistics} \\
单用户物理数据总量 & 1GB &
\texttt{information\_schema.tables.data\_length\ +\ index\_length} \\
单数据库物理数据量 & 512MB & \texttt{information\_schema.tables} \\
单表物理数据量 & 256MB & \texttt{information\_schema.tables} \\
单次分页返回行数 & 100 & API 参数 \\
单次查询最大返回行数 & 1000 & 查询服务层 \\
单次批量插入行数 & 500 & API 请求体 \\
单次批量修改影响行数 & 200 & 影响行数预估 \\
单次批量删除影响行数 & 200 & 影响行数预估 \\
查询最大 JOIN 数 & 8 & 查询 AST \\
查询最大子查询深度 & 3 & 查询 AST \\
查询最大执行时间 & 10 秒 & MySQL 会话级限制或应用层超时 \\
查询 AST 最大体积 & 64KB & 请求体大小 \\
头像文件大小 & 2MB & 上传服务 \\
背景图文件大小 & 8MB & 上传服务 \\
单用户资源总量 & 100MB & \texttt{user\_assets.file\_size} \\
\bottomrule
\end{longtable}

\hypertarget{data-limit-003-ux5bb9ux91cfux7edfux8ba1ux539fux5219}{%
\subsubsection{DATA-LIMIT-003
容量统计原则}\label{data-limit-003-ux5bb9ux91cfux7edfux8ba1ux539fux5219}}

\begin{itemize}
\tightlist
\item
  数据库、表、视图、字段、索引数量以 MySQL 元数据为准。
\item
  物理容量以 \texttt{information\_schema.tables} 的
  \texttt{data\_length\ +\ index\_length} 估算。
\item
  用户资源容量以系统库 \texttt{user\_assets.file\_size} 汇总。
\item
  容量限制校验必须发生在写入或创建操作之前。
\item
  统计值允许存在短暂延迟，但限制判断必须尽量使用实时查询。
\end{itemize}

\hypertarget{ux5b89ux5168ux8bbeux8ba1ux9644ux5f55}{%
\section{安全设计附录}\label{ux5b89ux5168ux8bbeux8ba1ux9644ux5f55}}

\hypertarget{ux5b89ux5168ux8bbeux8ba1}{%
\section{安全设计}\label{ux5b89ux5168ux8bbeux8ba1}}

\hypertarget{ux5b89ux5168ux76eeux6807}{%
\subsection{1. 安全目标}\label{ux5b89ux5168ux76eeux6807}}

系统必须满足以下安全目标：

\begin{itemize}
\tightlist
\item
  未登录用户不能访问工作台接口。
\item
  用户不能访问他人数据库。
\item
  前端不能直接执行任意 SQL。
\item
  动态 SQL 不产生注入风险。
\item
  Access Token 泄露后的有效时间受限。
\item
  Refresh Token 被服务端追踪和轮换。
\item
  高风险操作具备二次确认。
\item
  批量操作具备影响分析和容量限制。
\item
  数据库错误被精准归类并安全返回。
\item
  关键操作具备 traceId 和审计日志。
\end{itemize}

\hypertarget{ux53cc-token-ux5b89ux5168}{%
\subsection{2. 双 Token 安全}\label{ux53cc-token-ux5b89ux5168}}

\hypertarget{sec-auth-001-access-token}{%
\subsubsection{SEC-AUTH-001 Access
Token}\label{sec-auth-001-access-token}}

Access Token 用于访问普通 API。

\begin{longtable}[]{@{}ll@{}}
\toprule
属性 & 规则 \\
\midrule
\endhead
有效期 & 15 分钟 \\
签名密钥 & \texttt{JWT\_ACCESS\_SECRET} \\
传输方式 &
\texttt{Authorization:\ Bearer\ \textless{}accessToken\textgreater{}} \\
Payload & \texttt{userID}、\texttt{accountNo} \\
\bottomrule
\end{longtable}

Payload 禁止包含：

\begin{itemize}
\tightlist
\item
  密码。
\item
  密码哈希。
\item
  Refresh Token。
\item
  数据库列表。
\item
  用户资源路径。
\end{itemize}

\hypertarget{sec-auth-002-refresh-token}{%
\subsubsection{SEC-AUTH-002 Refresh
Token}\label{sec-auth-002-refresh-token}}

Refresh Token 用于刷新 Access Token。

\begin{longtable}[]{@{}ll@{}}
\toprule
属性 & 规则 \\
\midrule
\endhead
有效期 & 7 天 \\
签名密钥 & \texttt{JWT\_REFRESH\_SECRET} \\
存放位置 & HttpOnly Cookie \\
服务端存储 & token 哈希 \\
\bottomrule
\end{longtable}

Cookie 属性：

\begin{Shaded}
\begin{Highlighting}[]
\NormalTok{HttpOnly = true}
\NormalTok{Secure = true}
\NormalTok{SameSite = Strict}
\NormalTok{Path = /api/v1/auth}
\end{Highlighting}
\end{Shaded}

开发环境通过环境变量关闭 \texttt{Secure}，生产环境必须开启
\texttt{Secure}。

\hypertarget{sec-auth-003-refresh-token-rotation}{%
\subsubsection{SEC-AUTH-003 Refresh Token
Rotation}\label{sec-auth-003-refresh-token-rotation}}

刷新流程：

\begin{Shaded}
\begin{Highlighting}[]
\NormalTok{校验 Refresh Token}
\NormalTok{校验 token\_hash 存在且未撤销}
\NormalTok{撤销旧 Refresh Token}
\NormalTok{签发新 Access Token}
\NormalTok{签发新 Refresh Token}
\NormalTok{保存新 Refresh Token 哈希}
\NormalTok{设置新 Refresh Token Cookie}
\NormalTok{写入审计日志}
\end{Highlighting}
\end{Shaded}

被撤销的 Refresh Token 再次出现时，系统撤销同一 token family
下的全部活跃 Refresh Token。

\hypertarget{sec-auth-004-ux9000ux51faux767bux5f55}{%
\subsubsection{SEC-AUTH-004
退出登录}\label{sec-auth-004-ux9000ux51faux767bux5f55}}

退出登录流程：

\begin{Shaded}
\begin{Highlighting}[]
\NormalTok{校验 Refresh Token}
\NormalTok{撤销 Refresh Token}
\NormalTok{清除 Refresh Token Cookie}
\NormalTok{写入审计日志}
\end{Highlighting}
\end{Shaded}

\hypertarget{ux6388ux6743ux5b89ux5168}{%
\subsection{3. 授权安全}\label{ux6388ux6743ux5b89ux5168}}

\hypertarget{sec-authz-001-ux7528ux6237ux4e0aux4e0bux6587}{%
\subsubsection{SEC-AUTHZ-001
用户上下文}\label{sec-authz-001-ux7528ux6237ux4e0aux4e0bux6587}}

所有受保护接口必须经过鉴权中间件。鉴权中间件将当前用户 ID
写入请求上下文。

\hypertarget{sec-authz-002-ux6570ux636eux5e93ux5f52ux5c5eux6821ux9a8c}{%
\subsubsection{SEC-AUTHZ-002
数据库归属校验}\label{sec-authz-002-ux6570ux636eux5e93ux5f52ux5c5eux6821ux9a8c}}

所有数据库对象接口必须执行：

\begin{Shaded}
\begin{Highlighting}[]
\NormalTok{currentUserId + databaseId + status active}
\end{Highlighting}
\end{Shaded}

查询 \texttt{user\_databases} 成功后才能获得 \texttt{schema\_name}。

\hypertarget{sec-authz-003-ux5bf9ux8c61ux5f52ux5c5eux7ee7ux627f}{%
\subsubsection{SEC-AUTHZ-003
对象归属继承}\label{sec-authz-003-ux5bf9ux8c61ux5f52ux5c5eux7ee7ux627f}}

表、字段、索引、约束、视图、数据行、事务均继承数据库归属权限。没有通过数据库归属校验的请求不能继续访问
MySQL 用户物理库。

\hypertarget{sec-authz-004-ux4e8bux52a1ux5f52ux5c5eux6821ux9a8c}{%
\subsubsection{SEC-AUTHZ-004
事务归属校验}\label{sec-authz-004-ux4e8bux52a1ux5f52ux5c5eux6821ux9a8c}}

事务提交、回滚和事务内写操作必须校验：

\begin{Shaded}
\begin{Highlighting}[]
\NormalTok{transactionId + currentUserId + databaseId + status active}
\end{Highlighting}
\end{Shaded}

\hypertarget{sql-ux5b89ux5168}{%
\subsection{4. SQL 安全}\label{sql-ux5b89ux5168}}

\hypertarget{sec-sql-001-ux503cux53c2ux6570ux5316}{%
\subsubsection{SEC-SQL-001
值参数化}\label{sec-sql-001-ux503cux53c2ux6570ux5316}}

所有用户输入值必须使用参数化查询。

\begin{Shaded}
\begin{Highlighting}[]
\KeywordTok{WHERE}\NormalTok{ \textasciigrave{}name\textasciigrave{} }\OperatorTok{=}\NormalTok{ ?}
\end{Highlighting}
\end{Shaded}

\hypertarget{sec-sql-002-ux6807ux8bc6ux7b26ux6821ux9a8c}{%
\subsubsection{SEC-SQL-002
标识符校验}\label{sec-sql-002-ux6807ux8bc6ux7b26ux6821ux9a8c}}

数据库显示名、表名、字段名、索引名、约束名、视图名必须符合：

\begin{Shaded}
\begin{Highlighting}[]
\NormalTok{\^{}[A{-}Za{-}z][A{-}Za{-}z0{-}9\_]\{0,63\}$}
\end{Highlighting}
\end{Shaded}

\hypertarget{sec-sql-003-ux6807ux8bc6ux7b26ux6765ux6e90ux6821ux9a8c}{%
\subsubsection{SEC-SQL-003
标识符来源校验}\label{sec-sql-003-ux6807ux8bc6ux7b26ux6765ux6e90ux6821ux9a8c}}

动态标识符必须满足：

\begin{itemize}
\tightlist
\item
  数据库真实名来自 \texttt{user\_databases.schema\_name}。
\item
  表名来自 \texttt{information\_schema.tables}。
\item
  字段名来自 \texttt{information\_schema.columns}。
\item
  索引名来自 \texttt{information\_schema.statistics}。
\item
  约束名来自 \texttt{information\_schema.table\_constraints}。
\item
  视图名来自 \texttt{information\_schema.views}。
\end{itemize}

\hypertarget{sec-sql-004-ux7edfux4e00ux8f6cux4e49}{%
\subsubsection{SEC-SQL-004
统一转义}\label{sec-sql-004-ux7edfux4e00ux8f6cux4e49}}

后端通过统一工具封装反引号。业务代码不能手写反引号拼接。

\hypertarget{sec-sql-005-gui-ux67e5ux8be2ux9650ux5236}{%
\subsubsection{SEC-SQL-005 GUI
查询限制}\label{sec-sql-005-gui-ux67e5ux8be2ux9650ux5236}}

GUI 查询必须执行：

\begin{itemize}
\tightlist
\item
  主表存在校验。
\item
  JOIN 表存在校验。
\item
  派生表子查询校验。
\item
  \texttt{IN}、\texttt{EXISTS}、标量子查询校验。
\item
  字段存在校验。
\item
  操作符白名单校验。
\item
  聚合函数白名单校验。
\item
  分组字段校验。
\item
  排序字段校验。
\item
  分页上限校验。
\item
  JOIN 数量上限校验。
\item
  子查询深度上限校验。
\item
  查询 AST 体积上限校验。
\item
  查询超时控制。
\end{itemize}

GUI 查询提交的是结构化查询 AST，不直接提交原始 SQL 字符串。后端只生成
\texttt{SELECT} 语句，禁止在查询构造器中生成
\texttt{INSERT}、\texttt{UPDATE}、\texttt{DELETE}、\texttt{DROP}、\texttt{ALTER}
等写操作。

V6 已在服务端执行 AST
校验、用户数据库归属校验、字段作用域解析、标识符引用转义、参数绑定、JOIN
数量限制、子查询深度限制、分页限制和查询执行时间提示。前端高级 AST
面板也只提交结构化 JSON，不开放原始 SQL 编辑器。

\hypertarget{sec-sql-006-ux5199ux64cdux4f5cux9650ux5236}{%
\subsubsection{SEC-SQL-006
写操作限制}\label{sec-sql-006-ux5199ux64cdux4f5cux9650ux5236}}

修改和删除必须携带 GUI 条件对象。后端禁止执行无条件 \texttt{UPDATE}
和无条件 \texttt{DELETE}。

批量写操作必须执行影响行数预估。超过阈值时返回
\texttt{CONFIRMATION\_REQUIRED}。

\hypertarget{ux9ad8ux98ceux9669ux64cdux4f5cux786eux8ba4}{%
\subsection{5.
高风险操作确认}\label{ux9ad8ux98ceux9669ux64cdux4f5cux786eux8ba4}}

\hypertarget{sec-confirm-001-ux5f3aux5236ux786eux8ba4}{%
\subsubsection{SEC-CONFIRM-001
强制确认}\label{sec-confirm-001-ux5f3aux5236ux786eux8ba4}}

以下操作必须二次确认，用户偏好不能关闭：

\begin{itemize}
\tightlist
\item
  删除数据库。
\item
  删除表。
\item
  删除字段。
\item
  删除视图。
\item
  触发外键级联删除。
\item
  事务回滚。
\end{itemize}

确认方式：

\begin{itemize}
\tightlist
\item
  删除数据库输入数据库显示名。
\item
  删除表输入表名。
\item
  删除字段输入字段名。
\item
  删除视图输入视图名。
\item
  级联删除展示影响表和影响行数后确认。
\item
  事务回滚展示事务 ID 后确认。
\end{itemize}

\hypertarget{sec-confirm-002-ux504fux597dux786eux8ba4}{%
\subsubsection{SEC-CONFIRM-002
偏好确认}\label{sec-confirm-002-ux504fux597dux786eux8ba4}}

以下操作达到阈值时二次确认，用户偏好控制同类提示：

\begin{itemize}
\tightlist
\item
  批量插入。
\item
  批量修改。
\item
  批量删除。
\item
  批量结构变更。
\end{itemize}

用户在确认弹窗勾选``本类操作不再提示''后，系统更新
\texttt{user\_preferences} 中对应字段。

\hypertarget{sec-confirm-003-ux786eux8ba4ux5931ux8d25}{%
\subsubsection{SEC-CONFIRM-003
确认失败}\label{sec-confirm-003-ux786eux8ba4ux5931ux8d25}}

确认文本不匹配时返回：

\begin{Shaded}
\begin{Highlighting}[]
\NormalTok{HTTP 412}
\NormalTok{error.type = CONFIRMATION\_REQUIRED}
\NormalTok{error.code = CONFIRMATION\_TEXT\_MISMATCH}
\end{Highlighting}
\end{Shaded}

\hypertarget{ux4e8bux52a1ux5b89ux5168}{%
\subsection{6. 事务安全}\label{ux4e8bux52a1ux5b89ux5168}}

\hypertarget{sec-tx-001-ux4e8bux52a1ux8fdeux63a5}{%
\subsubsection{SEC-TX-001
事务连接}\label{sec-tx-001-ux4e8bux52a1ux8fdeux63a5}}

事务会话绑定后端 MySQL 连接。连接只服务创建该事务的用户。

\hypertarget{sec-tx-002-ux4e8bux52a1ux8d85ux65f6}{%
\subsubsection{SEC-TX-002
事务超时}\label{sec-tx-002-ux4e8bux52a1ux8d85ux65f6}}

事务有效期固定为 120 秒。超时后系统自动回滚并释放连接。

\hypertarget{sec-tx-003-ux9694ux79bbux7ea7ux522b}{%
\subsubsection{SEC-TX-003
隔离级别}\label{sec-tx-003-ux9694ux79bbux7ea7ux522b}}

支持的隔离级别：

\begin{Shaded}
\begin{Highlighting}[]
\NormalTok{READ COMMITTED}
\NormalTok{REPEATABLE READ}
\NormalTok{SERIALIZABLE}
\end{Highlighting}
\end{Shaded}

\hypertarget{sec-tx-004-ux5ba1ux8ba1}{%
\subsubsection{SEC-TX-004 审计}\label{sec-tx-004-ux5ba1ux8ba1}}

事务开始、提交、回滚、超时回滚和事务内写操作都写入审计日志。

\hypertarget{ux7cbeux786eux9519ux8befux5b89ux5168}{%
\subsection{7.
精确错误安全}\label{ux7cbeux786eux9519ux8befux5b89ux5168}}

\hypertarget{sec-error-001-ux9519ux8befux5f52ux7c7b}{%
\subsubsection{SEC-ERROR-001
错误归类}\label{sec-error-001-ux9519ux8befux5f52ux7c7b}}

后端必须区分：

\begin{itemize}
\tightlist
\item
  参数校验错误。
\item
  鉴权错误。
\item
  授权错误。
\item
  确认缺失。
\item
  容量超限。
\item
  MySQL 连接错误。
\item
  MySQL 约束错误。
\item
  MySQL 执行错误。
\item
  服务端内部错误。
\end{itemize}

\hypertarget{sec-error-002-mysql-ux9519ux8befux8fd4ux56de}{%
\subsubsection{SEC-ERROR-002 MySQL
错误返回}\label{sec-error-002-mysql-ux9519ux8befux8fd4ux56de}}

MySQL 错误响应包含：

\begin{itemize}
\tightlist
\item
  \texttt{errno}
\item
  \texttt{sqlState}
\item
  \texttt{sqlMessage}
\item
  \texttt{databaseName}
\item
  \texttt{tableName}
\item
  \texttt{columnName}
\item
  \texttt{constraintName}
\end{itemize}

响应禁止包含：

\begin{itemize}
\tightlist
\item
  数据库连接密码。
\item
  Access Token。
\item
  Refresh Token。
\item
  完整 SQL 中的敏感参数值。
\end{itemize}

\hypertarget{sec-error-003-ux6279ux91cfux64cdux4f5cux9519ux8bef}{%
\subsubsection{SEC-ERROR-003
批量操作错误}\label{sec-error-003-ux6279ux91cfux64cdux4f5cux9519ux8bef}}

批量操作失败时必须返回：

\begin{itemize}
\tightlist
\item
  失败行下标。
\item
  失败字段。
\item
  失败原因。
\item
  MySQL 错误号。
\item
  SQLSTATE。
\item
  traceId。
\end{itemize}

\hypertarget{ux6587ux4ef6ux8d44ux6e90ux5b89ux5168}{%
\subsection{8.
文件资源安全}\label{ux6587ux4ef6ux8d44ux6e90ux5b89ux5168}}

\hypertarget{sec-file-001-ux8d44ux6e90ux5b58ux50a8}{%
\subsubsection{SEC-FILE-001
资源存储}\label{sec-file-001-ux8d44ux6e90ux5b58ux50a8}}

资源实体保存到服务端文件目录。系统库保存资源元信息。

\hypertarget{sec-file-002-ux6587ux4ef6ux6821ux9a8c}{%
\subsubsection{SEC-FILE-002
文件校验}\label{sec-file-002-ux6587ux4ef6ux6821ux9a8c}}

上传资源必须校验：

\begin{itemize}
\tightlist
\item
  文件大小。
\item
  MIME 类型。
\item
  扩展名。
\item
  文件内容签名。
\item
  校验和。
\end{itemize}

允许的图片 MIME 类型：

\begin{Shaded}
\begin{Highlighting}[]
\NormalTok{image/png}
\NormalTok{image/jpeg}
\NormalTok{image/webp}
\end{Highlighting}
\end{Shaded}

SVG 文件固定拒绝。

\hypertarget{sec-file-003-ux6587ux4ef6ux540d}{%
\subsubsection{SEC-FILE-003
文件名}\label{sec-file-003-ux6587ux4ef6ux540d}}

服务端生成随机资源键。用户原始文件名只作为展示元信息保存。

\hypertarget{sec-file-004-ux8d44ux6e90ux5f52ux5c5e}{%
\subsubsection{SEC-FILE-004
资源归属}\label{sec-file-004-ux8d44ux6e90ux5f52ux5c5e}}

读取、替换、删除资源必须校验 \texttt{asset.user\_id\ =\ currentUserId}。

\hypertarget{ux5bb9ux91cfux4e0eux9891ux7387ux9650ux5236}{%
\subsection{9.
容量与频率限制}\label{ux5bb9ux91cfux4e0eux9891ux7387ux9650ux5236}}

\hypertarget{sec-limit-001-ux5bb9ux91cfux9650ux5236}{%
\subsubsection{SEC-LIMIT-001
容量限制}\label{sec-limit-001-ux5bb9ux91cfux9650ux5236}}

服务端配置固定限制。默认值如下：

\begin{longtable}[]{@{}lll@{}}
\toprule
限制项 & 默认值 & 安全目的 \\
\midrule
\endhead
单用户数据库数量 & 10 & 防止大量创建 schema \\
单数据库表数量 & 100 & 防止元数据膨胀 \\
单数据库视图数量 & 50 & 防止复杂视图滥用 \\
单表字段数量 & 80 & 防止超宽表 \\
单表索引数量 & 16 & 防止索引滥用拖慢写入 \\
单用户物理数据总量 & 1GB & 控制本地 MySQL 容量 \\
单数据库物理数据量 & 512MB & 控制单库容量 \\
单表物理数据量 & 256MB & 控制单表容量 \\
单次分页返回行数 & 100 & 防止接口返回过大 \\
单次查询最大返回行数 & 1000 & 防止大结果集压垮前端 \\
单次批量插入行数 & 500 & 控制批量写入压力 \\
单次批量修改影响行数 & 200 & 防止误更新大量数据 \\
单次批量删除影响行数 & 200 & 防止误删大量数据 \\
查询最大 JOIN 数 & 8 & 控制查询复杂度 \\
查询最大子查询深度 & 3 & 防止嵌套查询过深 \\
查询最大执行时间 & 10 秒 & 防止长查询占用连接 \\
查询 AST 最大体积 & 64KB & 防止请求体过大 \\
头像文件大小 & 2MB & 控制头像上传 \\
背景图文件大小 & 8MB & 控制背景上传 \\
单用户资源总量 & 100MB & 控制用户文件容量 \\
\bottomrule
\end{longtable}

超过限制时返回：

\begin{Shaded}
\begin{Highlighting}[]
\NormalTok{HTTP 413}
\NormalTok{error.type = LIMIT\_EXCEEDED}
\NormalTok{error.code = LIMIT\_EXCEEDED}
\end{Highlighting}
\end{Shaded}

错误详情必须包含超限项、当前值和最大值，不返回任何敏感路径或 SQL 文本。

\hypertarget{sec-limit-002-ux9891ux7387ux9650ux5236}{%
\subsubsection{SEC-LIMIT-002
频率限制}\label{sec-limit-002-ux9891ux7387ux9650ux5236}}

服务端限制：

\begin{itemize}
\tightlist
\item
  登录尝试频率。
\item
  Token 刷新频率。
\item
  创建数据库频率。
\item
  删除数据库频率。
\item
  批量写操作频率。
\item
  GUI 查询频率。
\end{itemize}

\hypertarget{sec-limit-003-ux767bux5f55ux5931ux8d25ux51b7ux5374}{%
\subsubsection{SEC-LIMIT-003
登录失败冷却}\label{sec-limit-003-ux767bux5f55ux5931ux8d25ux51b7ux5374}}

同一账号和同一 IP 多次登录失败后，系统启用冷却。冷却期间登录接口返回
\texttt{RATE\_LIMITED}。

\hypertarget{ux5ba1ux8ba1ux4e0eux8ffdux8e2a}{%
\subsection{10. 审计与追踪}\label{ux5ba1ux8ba1ux4e0eux8ffdux8e2a}}

\hypertarget{sec-audit-001-traceid}{%
\subsubsection{SEC-AUDIT-001 traceId}\label{sec-audit-001-traceid}}

每个请求必须生成 traceId。traceId 同时出现在：

\begin{itemize}
\tightlist
\item
  接口响应。
\item
  服务端日志。
\item
  审计日志。
\end{itemize}

\hypertarget{sec-audit-002-ux5ba1ux8ba1ux8303ux56f4}{%
\subsubsection{SEC-AUDIT-002
审计范围}\label{sec-audit-002-ux5ba1ux8ba1ux8303ux56f4}}

审计日志记录：

\begin{itemize}
\tightlist
\item
  注册。
\item
  登录。
\item
  Token 刷新。
\item
  退出登录。
\item
  偏好修改。
\item
  数据库创建、修改、删除。
\item
  表创建、修改、删除。
\item
  字段创建、修改、删除。
\item
  约束创建、删除。
\item
  索引创建、删除。
\item
  视图创建、修改、删除。
\item
  数据查询、新增、修改、删除。
\item
  GUI 查询。
\item
  事务开始、提交、回滚、超时回滚。
\item
  统计查看。
\item
  资源读取。
\end{itemize}

\hypertarget{sec-audit-003-ux5ba1ux8ba1ux5185ux5bb9}{%
\subsubsection{SEC-AUDIT-003
审计内容}\label{sec-audit-003-ux5ba1ux8ba1ux5185ux5bb9}}

审计日志包含：

\begin{itemize}
\tightlist
\item
  traceId。
\item
  userId。
\item
  databaseId。
\item
  transactionId。
\item
  objectType。
\item
  objectName。
\item
  actionType。
\item
  summary。
\item
  detailJson。
\item
  success。
\item
  errorCode。
\item
  errorMessage。
\item
  durationMs。
\item
  ipAddress。
\item
  userAgent。
\item
  createdAt。
\end{itemize}

审计日志禁止保存密码明文、Token 明文和完整敏感参数值。

\hypertarget{ux73afux5883ux53d8ux91cf}{%
\subsection{11. 环境变量}\label{ux73afux5883ux53d8ux91cf}}

服务启动必须校验以下环境变量：

\begin{Shaded}
\begin{Highlighting}[]
\NormalTok{JWT\_ACCESS\_SECRET}
\NormalTok{JWT\_REFRESH\_SECRET}
\NormalTok{DB\_HOST}
\NormalTok{DB\_USER}
\NormalTok{DB\_PASSWORD}
\NormalTok{DB\_NAME}
\NormalTok{PORT}
\end{Highlighting}
\end{Shaded}

缺少任一变量时服务启动失败。

代码中禁止提供生产可用默认密钥。

\hypertarget{ux524dux7aefux5b89ux5168}{%
\subsection{12. 前端安全}\label{ux524dux7aefux5b89ux5168}}

前端安全要求：

\begin{itemize}
\tightlist
\item
  不在前端保存 Refresh Token。
\item
  Access Token 只用于请求头，刷新失败立即清理登录态。
\item
  所有危险操作必须使用统一确认组件。
\item
  不在 URL 中携带 Token、密码、Refresh Token 或敏感查询参数。
\item
  表格行内编辑只进入草稿态，点击更新按钮才提交。
\item
  前端错误提示优先展示后端 \texttt{message}，但不能展示敏感字段。
\end{itemize}

\hypertarget{ux9632ux5fa1ux6027ux8bbeux8ba1ux6c47ux603b}{%
\subsection{13.
防御性设计汇总}\label{ux9632ux5fa1ux6027ux8bbeux8ba1ux6c47ux603b}}

防御性设计用于处理``用户误操作、恶意输入、服务异常、数据库异常、状态不一致''等情况。

\hypertarget{sec-def-001-ux9ed8ux8ba4ux62d2ux7edd}{%
\subsubsection{SEC-DEF-001
默认拒绝}\label{sec-def-001-ux9ed8ux8ba4ux62d2ux7edd}}

\begin{itemize}
\tightlist
\item
  未登录请求默认拒绝。
\item
  无数据库归属默认拒绝。
\item
  不符合标识符规则的名称默认拒绝。
\item
  超过容量限制默认拒绝。
\item
  查询 AST 超过复杂度限制默认拒绝。
\end{itemize}

\hypertarget{sec-def-002-ux5931ux8d25ux5b89ux5168}{%
\subsubsection{SEC-DEF-002
失败安全}\label{sec-def-002-ux5931ux8d25ux5b89ux5168}}

\begin{itemize}
\tightlist
\item
  创建数据库后系统表写入失败时，必须尝试删除已创建物理库。
\item
  事务超时必须自动回滚。
\item
  Refresh Token 复用疑似被盗时，撤销同一 token family。
\item
  批量修改和删除影响行数超过阈值时，不得直接执行。
\end{itemize}

\hypertarget{sec-def-003-ux6700ux5c0fux66b4ux9732}{%
\subsubsection{SEC-DEF-003
最小暴露}\label{sec-def-003-ux6700ux5c0fux66b4ux9732}}

\begin{itemize}
\tightlist
\item
  不返回真实物理数据库名。
\item
  不返回服务端文件绝对路径。
\item
  不返回密码、密码哈希、Token、Token 哈希。
\item
  MySQL 错误返回经过安全包装，只暴露可用于用户修正输入的信息。
\end{itemize}

\hypertarget{sec-def-004-ux53efux8ffdux8e2a}{%
\subsubsection{SEC-DEF-004
可追踪}\label{sec-def-004-ux53efux8ffdux8e2a}}

\begin{itemize}
\tightlist
\item
  所有关键失败都必须带 \texttt{traceId}。
\item
  所有写操作成功或失败都写入审计日志。
\item
  审计日志不记录敏感明文。
\end{itemize}

前端必须遵守：

\begin{itemize}
\tightlist
\item
  用户输入内容按文本渲染。
\item
  不渲染不可信 HTML。
\item
  Access Token 优先保存在内存。
\item
  页面刷新后通过 Refresh Token 恢复会话。
\item
  API 错误根据 \texttt{error.type} 和 \texttt{error.code} 精确展示。
\end{itemize}

\end{document}
